\documentclass[../Main.tex]{subfiles}


\begin{document}


\chapter{Martingale and Continuous Semimartingales Theory}

\section{Filtrations and Processes}

\subsection{Random Processes}

Throughout this section, let $(\Omega,\mathcal F,\PP)$ be a probability space, let $(E,d)$ be a metric space endowed with its Borel $\sigma$-field, and let the index set $T$ be either $\NN$ (discrete time) or an interval of $\RR$ (continuous time). A stochastic process is a family of $E$-valued random variables
\[ X=(X_t)_{t\in T}. \]

\defnr{Stochastic process}{stochastic-process}{ Let $(\Omega,\mathcal F,\PP)$ be a probability space, let $(E,d)$ be a metric space endowed with its Borel $\sigma$-field, and let $T$ be an index set, typically $T=\NN$ (discrete time) or $T\subseteq\RR$ (continuous time). An \emph{$E$-valued stochastic process} indexed by $T$ is a family of random variables \[ X=(X_t)_{t\in T}, \] where each $X_t:\Omega\to E$ is measurable. } 

\defnr{Sample paths}{sample-paths}{
Let $X=(X_t)_{t\in T}$ be an $E$-valued stochastic process. For every $\omega\in\Omega$, the mapping
\[ T\ni t\longmapsto X_t(\omega)\in E \]
is called the \emph{sample path} (or \emph{trajectory}) of $X$ associated with $\omega$.
}

\defnr{Graph of a stochastic process}{graph-process}{ The \emph{graph} of a sample path of a stochastic process $X=(X_t)_{t\in T}$ corresponding to $\omega\in\Omega$ is the subset of $T\times E$ defined by \[ \operatorname{Graph}(X(\omega)) =\{(t,X_t(\omega)):t\in T\}. \] When $E=\RR$, this is the usual graph of the function $t\mapsto X_t(\omega)$. }


\defnr{Modification}{modification}{
Let $X=(X_t)_{t\in T}$ and $\widetilde X=(\widetilde X_t)_{t\in T}$ be two $E$-valued stochastic processes defined on the same probability space and indexed by the same set $T$. We say that $\widetilde X$ is a \emph{modification} of $X$ if
\[
\forall t\in T,\qquad \PP(\widetilde X_t=X_t)=1.
\]
}

A modification preserves all finite-dimensional distributions. Consequently, any probabilistic property depending only on finite-dimensional marginals is shared by both processes. In contrast, pathwise properties may differ substantially. For example, when $T=\RR_+$ and $E=\RR$, a process may admit a modification with continuous sample paths even though its original sample paths are almost surely discontinuous.

\defnr{Indistinguishability}{def:indistinguishable}{
The processes $X$ and $\widetilde X$ are said to be \emph{indistinguishable} if there exists a measurable set $N\subseteq\Omega$ with $\PP(N)=0$ such that
\[
\forall \omega\in\Omega\setminus N,\ \forall t\in T,\qquad X_t(\omega)=\widetilde X_t(\omega).
\]
Equivalently,
\[
\PP\!\left(\forall t\in T,\ X_t=\widetilde X_t\right)=1,
\]
keeping in mind that, for an arbitrary index set $T$, the event above need not be measurable.
}

Every pair of indistinguishable processes are modifications of one another, but the converse is generally false. Indistinguishable processes coincide almost surely as functions of time, whereas modifications agree only at each fixed time. Throughout these notes, processes are identified up to indistinguishability. Thus, uniqueness statements should always be understood in this sense.

When $T=I\subseteq\RR$ is an interval, the distinction disappears under mild regularity assumptions. If both $X$ and $\widetilde X$ have continuous sample paths (or merely right-continuous or left-continuous sample paths) outside a $\PP$-null set, then
\[
X\text{ is a modification of }\widetilde X
\quad\Longleftrightarrow\quad
X\text{ and }\widetilde X\text{ are indistinguishable}.
\]
Indeed, equality holds almost surely at every rational time by countability, and continuity extends this equality to every $t\in I$.

\thmr{Kolmogorov's Continuity Theorem}{thm:kolmogorov-lemma}{
Let $I\subseteq\RR$ be a bounded interval, and let $X=(X_t)_{t\in I}$ be an $E$-valued stochastic process, where $(E,d)$ is complete. Suppose there exist constants $q,\varepsilon,C>0$ such that
\[
\EE[d(X_s,X_t)^q]\le C|t-s|^{1+\varepsilon},\qquad s,t\in I.
\]
Then $X$ admits a modification $\widetilde X$ whose sample paths are Hölder continuous of every order $\alpha\in(0,\varepsilon/q)$. More precisely, for every $\alpha\in(0,\varepsilon/q)$ and every $\omega\in\Omega$, there exists a finite constant $C_\alpha(\omega)$ such that
\[
d(\widetilde X_s(\omega),\widetilde X_t(\omega))
\le C_\alpha(\omega)|t-s|^\alpha,\qquad s,t\in I.
\]
In particular, $X$ admits a continuous modification, which is unique up to indistinguishability.
}

\noindent\textbf{Remarks.}
\begin{enumerate}
\item[(i)] If $I$ is unbounded, for example $I=\RR_+$, the theorem can be applied successively on bounded intervals $[0,1],[1,2],\ldots$. This yields a modification whose sample paths are locally Hölder continuous of every order $\alpha\in(0,\varepsilon/q)$.

\item[(ii)] It is sufficient to prove the theorem for one fixed exponent $\alpha\in(0,\varepsilon/q)$. Applying the result to a sequence $\alpha_n\uparrow\varepsilon/q$ gives modifications that are necessarily indistinguishable, and therefore coincide almost surely.
\end{enumerate}

\lemr{Dyadic Hölder Lemma}{lem:dyadic-holder}{
Let $D=\{i2^{-n}:n\ge1,\ 0\le i<2^n\}\subset[0,1)$, and let $f:D\to E$, where $(E,d)$ is a metric space. Assume that there exist constants $\alpha>0$ and $K<\infty$ such that
\[
d\!\left(f((i-1)2^{-n}),f(i2^{-n})\right)\le K2^{-n\alpha}
\]
for every $n\ge1$ and every $i=1,\ldots,2^n-1$. Then
\[
d(f(s),f(t))\le\frac{2K}{1-2^{-\alpha}}|t-s|^\alpha,\qquad s,t\in D.
\]
}


\subsection{Filtration} \label{sec:filtration}
\defnr{Filtration}{filtration}{A filtration on $(\Omega, \A, \PP)$ is an increasing family of sub-$\sigma$-fields of $\A$. The family in the most applications of this book is indexed by time, which is either discrete ($\NN$) or continuous ($\RR_+$). 
\begin{itemize}
    \item \textbf{The discrete-time case:} A discrete-time filtration is a family $(\F_n)_{n \in \NN}$ of $\sigma$-fields such that :
    \[ \F_0 \subset \F_1 \subset \F_2 \subset \cdots \subset \A \]
    \item \textbf{The continuous-time case:} A continuous-time filtration is a family $(\F_t)_{t \in \RR_+}$ of $\sigma$-fields such that :
    \[ \F_0 \subset \F_s \subset \F_t \subset \F \]
\end{itemize}

We also say that $(\Omega, \A, (\F_n)_{n \in \NN}, \PP)$ and   $(\Omega, \A, (\F_t)_{t\geq 0}, \PP)$  are filtered probability spaces.}

The $\sigma$-field $\F_n$ (resp. $\F_t$) gathers the information available at time $n$ (resp. $t$), events that are $\F_n$-measurable (resp. $\F_t$-measurable) are interpreted as those that depend only on what has happened up to time $n$ (resp. $t$). We will write : 
\[\F_\infty = 
\begin{cases}
\bigvee_{n \in \NN} \F_n := \sigma \left( \bigcup_{n \in \NN} \F_n \right) & \text{in the discrete-time case} \\
\bigvee_{t \geq 0} \F_t := \sigma \left( \bigcup_{t \geq 0} \F_t \right) & \text{in the continuous-time case}
\end{cases}\]

If $(X_n)_{n \in \NN}$ (resp. $(X_t)_{t \geq 0}$) is a random process, then, for every $n \in \NN$ (resp. $t \geq 0$), we may define $\F_n^X$ (resp. $\F_t^X$) as the smallest $\sigma$-field on $\Omega$ for which $X_0, X_1, \ldots, X_n$ (resp. $X_s \quad s\leq t$) are measurables. In other words, we have :
\[
\begin{cases}
    \F_n^X = \sigma(X_0, X_1, \ldots, X_n) & \text{in the discrete-time case} \\
    \F_t^X = \sigma(X_s : s \leq t) & \text{in the continuous-time case}
\end{cases}
\]
Then $(\F_n^X)_{n \in \NN}$ (resp. $(\F_t^X)_{t \geq 0}$) is a filtration called \textit{the canonical filtration} of $(X_n)_{n \in \NN}$ (resp. $(X_t)_{t \geq 0}$).

\textbf{Remark:}
Let $(\F_t)_{0 \le t \le \infty}$ be a filtration on $(\Omega, \F, P)$. We set, for every $t \ge 0$

$$\F_{t+} = \bigcap_{s > t} \F_s,$$

and $\F_{\infty+} = \F_\infty$. Note that $\F_t \subset \F_{t+}$ for every $t \in [0, \infty]$. The collection $(\F_{t+})_{0 \le t \le \infty}$ is also a filtration. We say that the filtration $(\F_t)$ is right-continuous if

$$\F_{t+} = \F_t, \qquad \forall t \ge 0.$$

Let $(\F_t)$ be a filtration and let $\N$ be the class of all $(\F_\infty, P)$-negligible sets (i.e. $A \in \N$ if there exists an $A' \in \F_\infty$ such $A \subset A'$ and $P(A') = 0$). The filtration is said to be complete if $\N \subset \F_0$ (and thus $\N \subset \F_t$ for every $t$).

If $(\F_t)$ is not complete, it can be completed by setting $\F'_t = \F_t \vee \sigma(\N)$, for every $t \in [0, \infty]$, using the notation $\F_t \vee \sigma(\N)$ for the smallest $\sigma$-field that contains both $\F_t$ and $\sigma(\N)$ (recall that $\sigma(\N)$ is the $\sigma$-field generated by $\N$). We will often apply this completion procedure to the canonical filtration of a random process $(X_t)_{t \ge 0}$, and call the resulting filtration the \textit{completed canonical filtration} of $X$.

\textbf{Assumption:}
\textit{In this book we will consider only filtrations that are right-continuous and complete.}

\subsection{Measurability of processes}
For the discrete-time case, we have only one definition for measurablity with respect to a filtration, but for the continuous-time case, we have 3 types/definitions of measurability. We start with the discrete-time case.
\defnr{Adapted Process [discrete]}{adapted-discrete}{
    A random process $(X_n)_{n \in \NN}$ is said to be adapted to the filtration $(\F_n)_{n \in \NN}$ if, for every $n \in \NN$, $X_n$ is $\F_n$-measurable.
}
In the continuous-time case, we have the following definitions.
\defnr{Measurable Process [continuous]}{measurable-continuous}{
    A process $X = (X_t)_{t \ge 0}$ with values in a measurable space $(E, \E)$ is said to be measurable if the mapping $$(\omega, t) \mapsto X_t(\omega)$$ defined on $\Omega \times \RR_+$ equipped with the product $\sigma$-field $\F \otimes \B(\RR_+)$ is measurable. (We recall that $\B(\RR_+)$ stands for the Borel $\sigma$-field of $\RR$.)
}
This is stronger than saying that, for every $t \ge 0$, $X_t$ is $\F$-measurable. On the other hand, considering for instance the case where $E = \RR$, it is easy to see that if the sample paths of $X$ are continuous, or only right-continuous, the fact that $X_t$ is $\F$-measurable for every $t$ implies that the process is measurable in the previous sense -- see the argument in the proof of Proposition \ref{prop:adapted-progressive-right-continuous} below.

\defnr{Adapted and Progressive Process}{adapted-continuous}{
A random process $(X_t)_{t \ge 0}$ with values in a measurable space $(E, \E)$ is called \textbf{adapted} if, for every $t \ge 0$, $X_t$ is $\F_t$-measurable. This process is said to be \textbf{progressive} if, for every $t \ge 0$, the mapping
$$(\omega, s) \mapsto X_s(\omega)$$
defined on $\Omega \times [0, t]$ is measurable for the $\sigma$-field $\F_t \otimes \B([0, t])$.
}

\proppr{adapted-progressive-right-continuous}{
Let $(X_t)_{t \ge 0}$ be a random process with values in a metric space $(E, d)$ (equipped with its Borel $\sigma$-field). Suppose that $X$ is adapted and that the sample paths of $X$ are right-continuous (i.e. for every $\omega \in \Omega$, $t \mapsto X_t(\omega)$ is right-continuous). Then $X$ is progressive. The same conclusion holds if one replaces right-continuous by left-continuous.
}{
We treat only the case of right-continuous sample paths, as the other case is similar. Fix $t > 0$. For every $n \ge 1$ and $s \in [0, t]$, define a random variable $X^n_s$ by setting

$$X^n_s = X_{kt/n} \quad \text{if } s \in \left[(k - 1)\frac{t}{n}, k\frac{t}{n}\right), k \in \{1, \ldots, n\},$$

and $X^n_t = X_t$. The right-continuity of sample paths ensures that, for every $s \in [0, t]$ and $\omega \in \Omega$,

\[ X_s(\omega) = \lim_{n \to \infty} X^n_s(\omega). \]

On the other hand, for every Borel subset $A$ of $E$,

\[
\begin{aligned}
\{(\omega, s) \in \Omega \times [0, t] : X^n_s(\omega) \in A\}
&= (\{X_t \in A\} \times \{t\}) \\
&\quad \cup \left( \bigcup_{k=1}^n \left( \{X_{kt/n} \in A\} \times \left[ \frac{(k - 1)t}{n}, \frac{kt}{n} \right) \right) \right)
\end{aligned}
\]

which belongs to the $\sigma$-field $\F_t \otimes \B([0, t])$. Hence, for every $n \ge 1$, the mapping $(\omega, s) \mapsto X^n_s(\omega)$, defined on $\Omega \times [0, t]$, is measurable for $\F_t \otimes \B([0, t])$. Since a pointwise limit of measurable functions is also measurable, the same measurability property holds for the mapping $(\omega, s) \mapsto X_s(\omega)$ defined on $\Omega \times [0, t]$. It follows that the process $X$ is progressive. 
}

\textbf{The progressive $\sigma$-field.} The collection $\mathcal P$ of all sets $A\in\mathcal F\otimes\mathcal B(\RR_+)$ such that the process
\[ X_t(\omega)=\mathbf{1}_A(\omega,t) \]
is progressive forms a $\sigma$-field on $\Omega\times\RR_+$, which is called the progressive $\sigma$-field. A subset $A$ of $\Omega\times\RR_+$ belongs to $\mathcal P$ if and only if, for every $t\geq 0$,
\[ A\cap(\Omega\times[0,t])\in\mathcal F_t\otimes\mathcal B([0,t]). \]

One then verifies that a process $X$ is progressive if and only if the mapping
\[ (\omega,t)\longmapsto X_t(\omega) \]
is measurable on $\Omega\times\RR_+$ equipped with the $\sigma$-field $\mathcal P$.


\section{Martingales and stopping times} \label{sec:martingale_stoppig_time}
\subsection{Martingales} 

\defnr{Martingale }{Martingales}{
An adapted real-valued process $(X_n)_{n \in \NN}$ (resp. $(X_t)_{t \ge 0}$) such that $X_n \in L^1$ for every $n\in\NN$ (resp. $X_t \in L^1$ for every $t \ge 0$) is called
\begin{itemize}
\item a \textbf{martingale} if : 
\[ \begin{cases}
    \forall n \in \NN \quad  X_n = \EE[X_{n+1}|F_n]   & \text{For the discret time} \\ 
    \forall s \geq t \ge 0 \quad X_t = \EE[X_{s}|F_t]  & \text{For the continuous time}
\end{cases}\]
\item a \textbf{supermartingale} if : 
\[ \begin{cases}
    \forall n \in \NN \quad X_n \geq \EE[X_{n+1}|F_n] & \text{For the discret time} \\ 
    \forall s \geq t \ge 0 \quad X_t \geq \EE[X_{s}|F_t]   & \text{For the continuous time}
\end{cases}\]

\item a \textbf{submartingale} if : 
\[ \begin{cases}
    \forall n \in \NN \quad X_n \leq \EE[X_{n+1}|F_n]  & \text{For the discret time} \\ 
    \forall s \geq t \ge 0 \quad X_t \leq \EE[X_{s}|F_t]   & \text{For the continuous time}
\end{cases}\]

\end{itemize}
}
If $(X_n)_{n \in \NN}$ (resp.  $(X_t)_{t \ge 0}$) is a submartingale, $(-X_n)_{n \in \NN}$ (resp. $(-X_t)_{t \ge 0}$) is a supermartingale. For this reason, some of the results below are stated for supermartingales only, but the analogous results for submartingales immediately follow.

\propr{convexity_martingal}{
Let $\varphi : \RR \longrightarrow \RR_+$ be a convex function, and let $(X_n)_{n \in \NN}$ (resp. $(X_t)_{t \ge 0}$) be an adapted process, such that $\EE[\varphi(X_n)] < \infty$ for every $n \in \NN$ (resp. $\EE[\varphi(X_t)] < \infty$ for every $t \ge 0$).
\begin{itemize}
\item  If $(X_n)_{n \in \NN}$ (resp. $(X_t)_{t \ge 0}$) is a martingale, then $(\varphi(X_n))_{n \in \NN}$ (resp. $(\varphi(X_t))_{t \ge 0}$) is a submartingale.
\item  If $(X_n)_{n \in \NN}$ (resp. $(X_t)_{t \ge 0}$) is a submartingale, and if we assume in addition that $\varphi$ is increasing, $(\varphi(X_n))_{n \in \NN}$ (resp. $(\varphi(X_t))_{t \ge 0}$) is a submartingale.
\end{itemize}
}

In particular, if $X_n$ (resp. $X_t$) is a martingale, $|X_n|$ (resp. $|X_t|$) is a submartingale (and so is $X_n^2$, provided that $\EE[X_n^2] < \infty$ for every $n$) and, if $X_n$ (resp. $X_t$) is a submartingale, $(X_n)^+$ (resp. $(X_t)^+$) is also a submartingale.

% \propr{sup_martingale_finite}{
%     Let $(X_t)_{t \geq 0}$ be a submartingale or a supermartingale. Then, for every $t > 0$, 
%     \[ \sup_{0 \leq s \leq t} E[|X_s|] < \infty. \]
% }

\propr{finite_quad_var}{
    Let $(M_n)_{n\in \NN}$ (resp.  $(M_t)_{t \geq 0}$) be a square integrable martingale ($\in L^2$). Then :
    \begin{itemize}
        \item \textbf{Discrete-time case:} If $n \geq m \geq 1$. Then,
        \[\EE\left[ \sum_{k=m}^{n-1} (M_{k+1} - M_{k})^2 \mid \F_m \right] = \EE[M_n^2 - M_m^2 \mid \F_m] = \EE[(M_n - M_m)^2 \mid \F_m].\]
        In particular, 
        \[ \EE\left[ \sum_{i=m}^{n-1} (M_{k+1} - M_{t_{k}})^2 \right] = \EE[M_n^2 - M_m^2] = \EE[(M_n - M_m)^2]. \]
        \item \textbf{Continuous-time case:} If $0 \leq s < t$ and let $s = t_0 < t_1 < \dots < t_p = t$ be a subdivision of the interval $[s, t]$. Then,  
        \[ \EE\left[ \sum_{i=1}^{p} (M_{t_i} - M_{t_{i-1}})^2 \mid \F_s \right] = \EE[M_t^2 - M_s^2 \mid \F_s] = \EE[(M_t - M_s)^2 \mid \F_s]. \]
        In particular, \[ \EE\left[ \sum_{i=1}^{p} (M_{t_i} - M_{t_{i-1}})^2 \right] = \EE[M_t^2 - M_s^2] = \EE[(M_t - M_s)^2].\]
    \end{itemize}
}

% For the discrete-time martingale we can define the following operation with predictable processes that conserve martingality. 

% \defnr{Predictable Process}{predictable_process}{A sequence $(H_n)_{n \in \NN}$ of real random variables is called predictable if $H_n$ is bounded and $\F_{n-1}$-measurable, for every $n \in \NN$.}


% \proppr{predictable_martingale}{Let $(X_n)_{n \in \NN}$ be an adapted process, and let $(H_n)_{n \in \NN}$ be a predictable sequence. We set $(H \bullet X)_0 = 0$ and, for every integer $n \ge 1$:
% \[(H \bullet X)_n = H_1(X_1 - X_0) + H_2(X_2 - X_1) + \cdots + H_n(X_n - X_{n-1})\]
% Then,
% \begin{itemize}
%     \item  If $(X_n)_{n \in \NN}$ is a martingale, $((H \bullet X)_n)_{n \in \NN}$ is also a martingale.
%     \item If $(X_n)_{n \in \NN}$ is a supermartingale (resp. a submartingale), and $H_n \ge 0$ for every $n \in \NN$, then $((H \bullet X)_n)_{n \in \NN}$ is a supermartingale (resp. a submartingale).
% \end{itemize}}{
% Let us prove first assertion. Since the random variables $H_n$ are bounded, one immediately checks that the random variables $(H \bullet X)_n$ are integrable. It is also straightforward to verify that the process $((H \bullet X)_n)_{n \in \NN}$ is adapted (all random variables entering the definition of $(H \bullet X)_n$ are $\F_n$-measurable). Then it suffices to verify that, for every $n \in \NN$,

% \[ \EE[(H \bullet X)_{n+1} - (H \bullet X)_n \mid \F_n] = 0.\]

% However, $(H \bullet X)_{n+1} - (H \bullet X)_n = H_{n+1}(X_{n+1} - X_n)$ and since $H_{n+1}$ is $\F_n$-measurable, We can pull the measurable term out of the expectation, that

% \[\begin{aligned} 
%     \EE[H_{n+1}(X_{n+1} - X_n) \mid \F_n] 
%         &= H_{n+1} \EE[X_{n+1} - X_n \mid \F_n] \\ 
%         &= H_{n+1}(\EE[X_{n+1} \mid \F_n] - X_n) = 0. 
% \end{aligned}\]

% The proof of the second assertion is similar. 
% }

\subsection{Stopping times}
\defnr{Stopping time}{stopping_time}{
A random variable $T : \Omega \longrightarrow \NN \cup \{+\infty\}$ (resp. $T : \Omega \longrightarrow [0, \infty]$) is a \textbf{stopping time} of the filtration  $(\F_n)$ (resp. $(\F_t)$) if $\{T \leq n\} \in \F_n \quad \forall n \in\NN$ (resp. $\{T \leq t\} \in \F_t \quad \forall t \geq 0$)  The $\sigma$-field of the past before $T$ is then defined by

\[ 
\begin{cases}
    \F_T = \{A \in \F_{\infty} : \forall n \in \NN, A \cap \{T \leq n\} \in \F_n\}.    & \text{For the discrete time} \\
    \F_T = \{A \in \F_{\infty} : \forall t \geq 0, A \cap \{T \leq t\} \in \F_t\}.    & \text{For the continuous time} \\
\end{cases}\]

The reader will verify that $\F_T$ is indeed a $\sigma$-field.}

In what follows, discrete (resp. continuous) "stopping time" will mean stopping time of the filtration $(\F_n)$ (resp. $(\F_t)$) unless otherwise specified. If $T$ is a continuous stopping time, we also have $\{T < t\} \in \F_t$ for every $t > 0$, by considering $\{T < t\} = \bigcup_{n=1}^{\infty} \left\{ T \leq t - \frac{1}{n} \right\}$ , and moreover

\[ \{T = \infty\} = \left( \bigcup_{n=1}^{\infty} \{T \leq n\} \right)^c \in \F_{\infty}. \]

A good advantage of assuming the right continuity of the filtration is that, in this case, the definition of stopping times in the continuous time  case can be equivalently stated by replacing $\{T \leq t\} \in \F_t$ by $\{T < t\} \in \F_t$ for every $t \geq 0$. In fact : 
\[ \{T \leq t\} = \bigcap_{n=1}^{\infty} \left\{T < t + \frac{1}{n}\right\}. \in \F_{t+} = \F_t \]

\proppr{two_stopping_times}{
Let $S$ and $T$ be two stopping times
\begin{itemize} 
    \item If $S \leq T$. Then, $\F_S \subset \F_T$.
    \item We have in general, $S \wedge T$ and $S \vee T$ are also stopping times. Moreover $\F_{S \wedge T} = \F_S \cap \F_T$.
\end{itemize}
}{
    Can be easily proven by using the definition of stopping times and \sfield of stopping times and manipulation of sets.
}

\lempr{Measurability w.r.t. stopped filtration }{measurability_wrt_stopped_filtration}{
    Let $T$ be a stopping time. A function $\omega \mapsto Y(\omega)$ defined on the set $\{T < \infty\}$ and taking values in the measurable set $(E, \E)$ is $\F_T$-measurable if and only if, for every $n\in\NN$ (resp. $t \geq 0$), the restriction of $Y$ to the set $\{T = n\}$ is $\F_n$ measurable (resp. $\{T \leq t\}$ $\F_t$-measurable).
}{
For the discrete-time case, first assume that, for every $n \in\NN$, the restriction of $Y$ to $\{T =n \}$ is $\F_n$-measurable.
Then, for every measurable subset $A$ of $E$,
\[ \{Y \in A\} \cap \{T \leq  n \} = \bigcup_{k=0}^n \left( \{Y \in A\} \cap \{T = k\} \right) \in \F_n. \]
Letting $n \to \infty$, we first obtain that $\{Y \in A\} \in \F_\infty$, and then we deduce from the previous display that $\{Y \in A\} \in \F_T$.
Conversely, if $Y$ is $\F_T$-measurable, $\{Y \in A\} \in \F_T$ and thus $\{Y \in A\} \cap \{T = n\} =  \{Y \in A\} \cap \{T \leq n\} \setminus   \{Y \in A\} \cap \{T \leq n-1\} \in \F_n$, giving the desired result.

For the continuous-time case, first assume that, for every $t \ge 0$, the restriction of $Y$ to $\{T \le t\}$ is $\F_t$-measurable. Then, for every measurable subset $A$ of $E$,
\[ \{Y \in A\} \cap \{T \le t\} \in \F_t. \]
Letting $t \to \infty$, we first obtain that $\{Y \in A\} \in \F_\infty$, and then we deduce from the previous display that $\{Y \in A\} \in \F_T$.

Conversely, if $Y$ is $\F_T$-measurable, $\{Y \in A\} \in \F_T$ and thus $\{Y \in A\} \cap \{T \le t\} \in \F_t$, giving the desired result. 
}

\thmpr{Stopped process Measurability}{stopped_process_measurability}{
    Let $(Y_n)_{n \in \NN}$  (resp. $(Y_t)_{t \geq 0}$) be an adapted (resp. a progressive) process, and let $T$ be a stopping time. Then the random variable $Y_T$, which is defined on the event $\{T < \infty\}$ by $Y_T(\omega) = Y_{T(\omega)}(\omega)$, is $\F_T$-measurable.
}{
    Let $B$ be a Borel subset of $\RR$. Let's begin with the dicrete-time case, consider the restriction of $Y_T$ to the set $\{T = n\}$, which is given by $Y_n$ on this set. Since $Y_n$ is $\F_n$-measurable, the restriction of $Y_T$ to $\{T = n\}$ is also $\F_n$-measurable. Hence, by Lemma \ref{lem:measurability_wrt_stopped_filtration}, $Y_T$ is $\F_T$-measurable.
    For the continuous-time case. Let $t \ge 0$. The restriction to $\{T \le t\}$ of the function $\omega \mapsto X_T(\omega)$ is the composition of the two mappings
    \[
    \begin{aligned}
    \{T \le t\} \ni \omega &\mapsto (\omega, T(\omega) \wedge t) \\
    \F_t &\quad \F_t \otimes \B([0, t])
    \end{aligned}
    \]
    and
    \[
    \begin{aligned}
    \Omega \times [0, t] \ni (\omega, s) &\mapsto X_s(\omega) \\
    \F_t \otimes \B([0, t]) &\quad \E
    \end{aligned}
    \]
    which are both measurable (the first one since $T \wedge t$ is $\F_t$-measurable, by Proposition \ref{prop:two_stopping_times}, and the second one by the definition of a progressive process). It follows that the restriction to $\{T \le t\}$ of the function $\omega \mapsto X_T(\omega)$ is $\F_t$-measurable, which gives the desired result by the lemma \ref{lem:measurability_wrt_stopped_filtration}. 
}
If $T$ is a stopping time, then, for every $n \in \NN$ (resp. $t \geq 0$), $n \wedge T$ (resp. $t \wedge T$) is also a stopping time (Proposition \ref{prop:two_stopping_times}) and it follows from the preceding theorem that $Y_{n \wedge T}$ (resp. $Y_{t \wedge T}$) is $\F_{n \wedge T}$-measurable (resp. $\F_{t \wedge T}$-measurable), hence also $\F_n$-measurable (resp. $\F_t$-measurable) by Proposition \ref{prop:two_stopping_times}.

\proppr{sequence_stopping_time}{
    If $(S_n)$ is a sequence of stopping times then :
    \begin{itemize}
        \item if $(S_n)$ is monotone increasing, then $S=\lim \uparrow S_n$ is a stopping time. 
        \item If $(S_n)$ is monotone decreasing, then $S=\lim \downarrow S_n$ is a stopping time and
            \[\F_{S}=\bigcap_n \F_{S_n} \]
    \end{itemize}
}{
    \begin{itemize}
        \item If $(S_n)$ is monotone increasing, then for every $t \geq 0$, $\{S \leq t\} = \bigcap_{n=1}^{\infty} \{S_n \leq t\}$, which belongs to the \sfield $\F_t$ since $S_n$ is a stopping time for every $n$. Hence, $S$ is a stopping time. same argument for the discrete time case.
        \item  If $(S_n)$ is monotone decreasing, then for every $t \geq 0$, $\{S < t\} = \bigcup_{n=1}^{\infty} \{S_n < t\}$, which belongs to the \sfield $\F_t$ since $S_n$ is a stopping time for every $n$. Hence, $S$ is a stopping time by the fact that $\F_t$ is right-continuous [check remarks after the definition \ref{defn:stopping_time}]. 
        \item For the last assertion we first verify that $\F_S \subset \bigcap_n \F_{S_n}$ thanks to last proposition \ref{prop:two_stopping_times}. Conversely, if $A \in \bigcap_n \F_{S_n}$, then 
        \[ A \cap \{S < t\} = \bigcup_{n=1}^{\infty} A \cap \{S_n < t\} \in \F_t, \]
        which gives $A \in \F_S$ again this is allowed thanks to the right-continuity of the filtration.
    \end{itemize}
}

\proppr{majorant_is_stoping_time}{
    Let $T$ be a stopping time and let $S$ be an $\F_T$-measurable random variable with values in $[0, \infty]$, such that $S \geq T$. 
    Then $S$ is also a stopping time. In particular, if $T$ is a continuous stopping time,
    \[    T_n=  \frac{\lceil 2^n T \rceil}{2^n}\cdot \1_{\{T<\infty\}}+\infty \cdot \1_{\{T=\infty\}}, \quad n=0,1,2, \ldots \]
    defines a sequence of stopping times that decreases to $T$.
}{
    For the first assertion, we write, for every $n \in \NN$,
    \[ \{S \leq n\} = \{S \leq n\} \cap \{T \leq n \} \in \F_n \]
    since $\{S \leq n\}$ is $\F_T$-measurable. And similar argument holds for the continuous time case.  Further more, the second assertion follows since $T_n \ge T$, and $T_n$ is a function of $T$, hence $\F_T$-measurable, and $T_n \downarrow T$ as $n \uparrow \infty$ by construction.
    
}

Finally lets defining hitting time. 
\defnr{Hitting time}{hitting_time}{
    Let $(X_t)_{t \geq 0}$ be an adapted process with values in a metric space $(E, d)$. 
    \begin{enumerate}
        \item Assume that the sample paths of $X$ are right-continuous, and let $O$ be an open subset of $E$. Then 
        \[ T_O = \inf\{t \geq 0 : X_t \in O\} \] 
        is a stopping time of the filtration $(\F_{t})$.
        \item Assume that the sample paths of $X$ are continuous, and let $F$ be a closed subset of $E$. Then
        \[ T_F = \inf\{t \geq 0 : X_t \in F\} \] 
        is a stopping time (of the filtration $(\F_t)$).
    \end{enumerate}
}
In fact the stopping time property can be proven by considering the two sets :
\[ 
\begin{cases}
    \{T_O < t\} &= \displaystyle \bigcup_{s \in [0, t) \cap \QQ} \{X_s \in O\} \in \F_t, \\
    \{T_F \leq t\} &= \displaystyle \left\{ \inf_{0 \leq s \leq t} d(X_s, F) = 0 \right\} = \left\{ \inf_{s \in [0, t] \cap \QQ} d(X_s, F) = 0 \right\} \in \F_t.
\end{cases}
\]

\subsection{Two Important Lemmas}

\lempr{Optional stopping theorem, first version}{optional_stopping_theorem_v1}{
    Suppose that $(X_n)_{n \in \NN}$ is a martingale (resp. a supermartingale) and let $T$ be a stopping time. Then $(X_{n \wedge T})_{n \in \NN}$ is also a martingale (resp. a supermartingale). In particular, if the stopping time $T$ is bounded, one has $X_T \in L^1$, and
    \begin{align*}
    \EE[X_T] = \EE[X_0] \quad (\text{resp. } \EE[X_T] \leq \EE[X_0]).
    \end{align*}
}{
    Fix $n \in \NN$. We wish to show that $$\EE[X_{(n+1)\wedge T} \mid \F_n] = X_{n \wedge T}.$$
    Observe the pointwise decomposition 
    \[ X_{(n+1)\wedge T} - X_{n \wedge T} = (X_{n+1} - X_n)\1_{{T > n}},\] 
    Therefore it suffices to show 
    \[ \EE\bigl[(X_{n+1} - X_n),\1_{{T > n}} \mid \F_n\bigr] = 0.\]
    Since $T$ is a stopping time, the event $\{T > n\} = \{T \leq n\}^c$ belongs to $\F_n$, so the indicator $\1_{\{T > n\}}$ is $\F_n$-measurable. Taking it outside the conditional expectation and using the martingale property of $(X_n)$: 
    \begin{align*}
    \EE\bigl[(X_{n+1} - X_n),\1_{\{T > n\}} \mid \F_n\bigr] 
        &= \1_{\{T > n\}},\EE[X_{n+1} - X_n \mid \F_n] \\ 
        &= \1_{\{T > n\}} \cdot 0 \\
        &= 0.
    \end{align*}
    Hence $(X_{n \wedge T})_{n \in \NN}$ is a martingale.

    If the stopping time $T$ is bounded above by the constant $N$, $X_T = X_{N \wedge T}$ is in $L^1$, and $\EE[X_T] = \EE[X_{N \wedge T}] = \EE[X_0]$.
    For the supermartingale case just replace equalities with inequalities.
}



\lempr{}{increasing_stopping_time_submartingale}{
    Let $(X_n)_{n \in \NN}$ be a supermartingale, and let $S$ and $T$ be two bounded stopping times such that $S \le T$. Then
    \[ \EE[X_T] \le \EE[X_S]. \]
}{
    We began with the discrete-time case. Let $(X_n,\F_n)_{n\ge 0}$ be a supermartingale and let $S \le T \le N$ be bounded stopping times.
    Define the stopped process increments:
    \[ X_T - X_S = \sum_{k=1}^{N} \1_{{S<k\le T}}(X_k - X_{k-1}).\]
    Take expectation:
    \[ \EE[X_T - X_S] = \sum_{k=1}^{N}
    \EE\left[\1_{{S<k\le T}}(X_k - X_{k-1})\right]. \] 

    Since $S,T$ are stopping times, the event $\{S<k\le T\} \in \F_{k-1}$. Hence the indicator is $\F_{k-1}$-measurable. Using the supermartingale property, $ \EE[X_k - X_{k-1} \mid \F_{k-1}] \le 0$, we get
    \[ \EE!\left[\1_{{S<k\le T}}(X_k - X_{k-1})\right] \le 0.\]
    Summing over $k$, and using the fact that $X_T, X_S \in L^1$ (by Lemma \ref{lem:optional_stopping_theorem_v1}),
    \[ \EE[X_T - X_S] \le 0, \quad\Rightarrow\quad \EE[X_T] \le \EE[X_S]. \]
}


\section{Martingale convergence theory}

\subsection{Upcrossing numbers}
\subsubsection{Discret time}
Consider first a deterministic real sequence $y = (y_n)_{n \in \NN}$, and $a < b$, the \textbf{upcrossing number} of this sequence along $[a, b]$ before time $n$, denoted by $M_{ab}^y(n)$, is the largest integer $k$ such that there exists a strictly increasing finite sequence
    \[ m_1 < n_1 < m_2 < n_2 < \dots < m_k < n_k \]
of nonnegative integers smaller than or equal to $n$ with the properties $y_{m_i} \le a$ and $y_{n_i} \ge b$, for every $i \in \{1, \dots, k\}$. In what follows we consider a sequence $Y = (Y_n)_{n \in \NN}$ of real random variables, and the associated upcrossing number $M_{ab}^Y(n)$ is then an integer-valued random variable. where $M_{ab}^y(0) = 0$ if such a sequence does not exist.
We will use the following simple analytic lemma, whose proof is left to the reader.


\lemr{}{convergence_and_upcrossing_number_discrete}{ 
    The sequence $(y_n)_{n \in \NN}$ converges in $\overline{\RR}$ if and only if, for every choice of the rationals $a$ and $b$ with $a < b$, we have $M_{ab}^y(n) < \infty$.
    }
\lemr{Doob's Upcrossing Inequality}{Doob_upcrossing_inequality_discrete}{ 
    Let $X = (X_n)_{n \in \NN}$ be a supermartingale. Then, for every reals $a < b$ and every $n \in \NN$,
    $$E[M_{ab}^Y(n)] \leq \frac{1}{b-a}E[(Y_n-a)^-]$$
    }

\subsubsection{Continuous time}
Let $f : I \longrightarrow \RR$ be a function defined on a subset $I$ of $\RR_+$. If $a < b$, the upcrossing number of $f$ along $[a, b]$, denoted by $M_{ab}^f(I)$, is the maximal integer $k \geq 1$ such that there exists a finite increasing sequence $s_1 < t_1 < \dots < s_k < t_k$ of elements of $I$ such that $f(s_i) \leq a$ and $f(t_i) \geq b$ for every $i \in \{1, \dots, k\}$ (if, even for $k = 1$, there is no such subsequence, we take $M_{ab}^f(I) = 0$, and if such a subsequence exists for every $k \geq 1$, we take $M_{ab}^f(I) = \infty$). Upcrossing numbers are a convenient tool to study the regularity of functions.

In the next lemma, the notation

$$\lim_{s \downarrow \downarrow t} f(s) \quad (\text{resp. } \lim_{s \uparrow \uparrow t} f(s) )$$
means
$$\lim_{s \downarrow t, s > t} f(s) \quad (\text{resp. } \lim_{s \uparrow t, s < t} f(s) ).$$

We say that $g : \RR_+ \to \RR$ is \textbf{càdlàg} (for the French "continue à droite avec des limites à gauche") if $g$ is right-continuous and has left-limits at every $t > 0$.

\lemr{}{limits_and_upcrossing_number_continuous}{ 
    Let $D$ be a countable dense subset of $\RR_+$ and let $f$ be a real function defined on $D$. We assume that, for every $T \in D$,
\begin{itemize}
    \item  the function $f$ is bounded on $D \cap [0, T]$;
    \item  for all rationals $a$ and $b$ such that $a < b$, 
        $$M_{ab}^f(D \cap [0, T]) < \infty.$$
\end{itemize}
    Then, the right-limit
    $$f(t+) := \lim_{s \downarrow \downarrow t, s \in D} f(s)$$
    exists for every real $t \geq 0$, and similarly the left-limit
    $$f(t-) := \lim_{s \uparrow \uparrow t, s \in D} f(s)$$
    exists for every real $t > 0$. Furthermore, the function $g : \RR_+ \to \RR$ defined by $g(t) = f(t+)$ is càdlàg.
}
We omit the proof of this analytic lemma. It is important to note that the right and left-limits $f(t+)$ and $f(t-)$ are defined for every $t \geq 0$ ($t > 0$ in the case of $f(t-)$) and not only for $t \in D$.

\lemr{Doob's Upcrossing Inequality}{Doob_upcrossing_inequality_continuous}{ 
    Let $D$ be a countable dense subset of $\RR_+$  and let $X = (X_t)_{t \geq 0}$ be a supermartingale. Then, for every reals $a < b$ and every $t \geq 0$,
    $$\EE[M_{ab}^X(D \cap [0, T])] \leq \frac{1}{b - a} \EE[(X_T - a)^-]$$
}{
    Can be obtained by monotone convergence applied on the discret version Lemma \ref{lem:Doob_upcrossing_inequality_discrete} by considering a sequence of finite subsets of $D$.
}


\subsection{Lp Convergence and Doob's inequalities}

\thmpr{Doob's supermartingale inequalities}{doob_supermartingale_inequalities}{
\begin{enumerate}
    \item (Maximal inequality for submartingales) Let $(X_n)_{n\in\NN}$ (resp. $(X_t)_{t \geq 0}$) be a submartingale (with right-continuous sample paths in the continuous case). Then, for every $n\in\NN$ (resp. $t > 0$) and every $\lambda > 0$,
        \[
        \begin{cases}
            \displaystyle \lambda \PP\left( \sup_{0 \leq m \leq n} |X_m| > \lambda \right) \leq \EE[X_n^+] & \text{For the discrete case} \\
            \displaystyle \lambda \PP\left( \sup_{0 \leq s \leq t} |X_s| > \lambda \right) \leq \EE[X_t^+] & \text{For the continuous case} \\
        \end{cases}
        \]
    \item (Maximal inequality for supermartingales) Let $(X_n)_{n\in\NN}$ (resp. $(X_t)_{t \geq 0}$) be a supermartingale (with right-continuous sample paths in the continuous case). Then, for every $n\in\NN$ (resp. $t > 0$) and every $\lambda > 0$,
        \[
        \begin{cases}
            \displaystyle \lambda \PP\left( \sup_{0 \leq m \leq n} |X_m| > \lambda \right) \leq \EE[X_0] + \EE[X_n^-] & \text{For the discrete case} \\
            \displaystyle \lambda \PP\left( \sup_{0 \leq s \leq t} |X_s| > \lambda \right) \leq \EE[X_0] + \EE[X_t^-] & \text{For the continuous case} \\
        \end{cases}
        \]
    \item (Doob's inequality in $L^p$ for submartingales) Let  $(X_n)_{n\in\NN}$ (resp. $(X_t)_{t \geq 0}$) be a nonnegative submartingale (with right-continuous sample paths in the continuous case). Then, for every $n\in\NN$ (resp. $t > 0$) and every $p > 1$,
    \[
    \begin{cases}
        \displaystyle \EE\left[ \sup_{0 \leq m \leq n} X_m^p \right] \leq \left( \frac{p}{p - 1} \right)^p \EE[|X_n|^p] & \text{For the discrete case} \\
        \displaystyle \EE\left[ \sup_{0 \leq s \leq t} X_s^p \right] \leq \left( \frac{p}{p - 1} \right)^p \EE[|X_t|^p] & \text{For the continuous case} \\
    \end{cases}
    \]
    \item (Doob's inequality in $L^p$ for martingales) Let  $(X_n)_{n\in\NN}$ (resp. $(X_t)_{t \geq 0}$) be a martingale (with right-continuous sample paths in the continuous case). Then, for every $n\in\NN$ (resp. $t > 0$) and every $p > 1$,
    \[
    \begin{cases}
        \displaystyle \EE\left[ \sup_{0 \leq m \leq n} |X_m|^p \right] \leq \left( \frac{p}{p - 1} \right)^p \EE[|X_n|^p] & \text{For the discrete case} \\
        \displaystyle \EE\left[ \sup_{0 \leq s \leq t} |X_s|^p \right] \leq \left( \frac{p}{p - 1} \right)^p \EE[|X_t|^p] & \text{For the continuous case} \\
    \end{cases}
    \]
\end{enumerate}
Note that part (3) and (4) of the inequality is useful only if $\EE[|X_n|^p] < \infty$(resp. $\EE[|X_t|^p] < \infty$).
}{
\begin{itemize}
    \item Let us prove the first assertion. Consider the stopping time
    \[T = \inf\{k \geq 0 : X_k \geq a\}.\]
    Then, if
    \[ A = \left\{ \sup_{0 \leq k \leq n} X_k \geq a \right\}, \]
    we have $A = \{T \leq n\}$. On the other hand, by applying the preceding lemma \ref{lem:increasing_stopping_time_submartingale} to the stopping times $T \wedge n$ and $n$, we have
    \[ \EE[X_{T \wedge n}] \leq \EE[X_n]. \]
    Furthermore,
    \[ X_{T \wedge n} \geq a \1_A + X_n \1_{A^c}. \]
    By combining these two inequalities, we get
    \[ \EE[X_n] \geq a \PP(A) + \EE[X_n \1_{A^c}] \]
    and it follows that $a \PP(A) \leq \EE[X_n \1_A]$, which gives the first assertion of the theorem in the discrete case.   

    For the continuous case, Fix $t > 0$ and consider a countable dense subset $D$ of $\RR_+$ such that $0 \in D$ and $t \in D$. Then $D \cap [0, t]$ is the increasing union of a sequence $(D_m)_{m \geq 1}$ of finite subsets $[0, t]$ of the form $D_m = \{t^m_0, t^m_1, \dots, t^m_m\}$ where $0 = t^m_0 < t^m_1 < \dots < t^m_m = t$. For every fixed $m$, we can apply the discrete time case to the sequence $Y_n = X_{t^m_{n \wedge m}}$, which is a discrete submartingale with respect to the filtration $\G_n = \F_{t^m_{n \wedge m}}$. We get
    \[ \lambda\PP\left( \sup_{s \in D_m} |X_s| > \lambda \right) \leq  \EE[X_t^+].\] 

    Then, we observe that

    \[P\left( \sup_{s \in D_m} |X_s| > \lambda \right) \uparrow\PP\left( \sup_{s \in D \cap [0,t]} |X_s| > \lambda \right) \]
    when $m \uparrow \infty$. We have thus
    \[ \lambda\PP\left( \sup_{s \in D \cap [0,t]} |X_s| > \lambda \right) \leq  \EE[X_t^+]. \]
    Finally, the right-continuity of sample paths (and the fact that $t \in D$) ensures that
    \[ \sup_{s \in D \cap [0,t]} |X_s| = \sup_{s \in [0,t]} |X_s|. \]

    \item The proof of the (2) is very similar to the proof of the (1). We begin with the discrete case and introduce the stopping time 
    \[ R = \inf\{k \geq 0 : X_k \geq a\},\]
    and the event $B = \{R \leq n\}$. We have then $\EE[X_{R \wedge n}] \leq \EE[X_0]$, and $\EE[X_{R \wedge n}] \geq a \PP(B) + \EE[X_n \1_{B^c}]$. It follows that $a \PP(B) \leq \EE[X_0] - \EE[X_n \1_{B^c}]$, which leads to the desired result. Then for the continuous case, we can use the same argument as in assertion (1) by considering a countable dense subset $D$ of $\RR_+$ such that $0 \in D$ and $t \in D$, and then applying the discrete case to the sequence $Y_n = X_{t^m_{n \wedge m}}$ for every fixed $m$. 

    \item In order to prove the assertion (3) in the discrete case, we may assume that $\EE[(X_n)^p] < \infty$. Then, by Jensen's inequality for conditional expectations, we have, for every $0 \leq k \leq n$,
    \[ \EE[(X_k)^p] \leq \EE[\EE[X_n | \F_k]^p] \leq \EE[\EE[(X_n)^p | \F_n]] = \EE[(X_n)^p]. \]
    Denote $\widetilde{X}_n= \sup_{0 \leq m \leq n} X_m$. Since $\widetilde{X}_n \leq X_0 + \dots + X_n$, it follows that we have also $\EE[(\widetilde{X}_n)^p] < \infty$. By assertion (1) of the theorem, for every $a > 0$,
    \[ a \, \PP(\widetilde{X}_n \geq a) \leq \EE[X_n \1_{\{\widetilde{X}_n \geq a\}}]. \]

    We multiply each side of this inequality by $a^{p-2}$ and integrate with respect to Lebesgue measure on $(0, \infty)$. For the left side, we get
    \[ \int_{0}^{\infty} a^{p-1} \PP(\widetilde{X}_n \geq a) \, da = \EE \left[ \int_{0}^{\widetilde{X}_n} a^{p-1} da \right] = \frac{1}{p} \EE[(\widetilde{X}_n)^p] \]

    using the Fubini theorem. Similarly, we get for the right side
    \begin{align*}
        \int_{0}^{\infty} a^{p-2} \EE[X_n \1_{\{\widetilde{X}_n \geq a\}}] \, da &= \EE \left[ X_n \int_{0}^{\widetilde{X}_n} a^{p-2} da \right] \\
        &= \frac{1}{p-1} \EE[X_n (\widetilde{X}_n)^{p-1}] \\
        &\leq \frac{1}{p-1} \EE[(X_n)^p]^{\frac{1}{p}} \EE[(\widetilde{X}_n)^p]^{\frac{p-1}{p}}
    \end{align*}
    by the Hölder inequality. We thus get
    \[ \frac{1}{p} \EE[(\widetilde{X}_n)^p] \leq \frac{1}{p-1} \EE[(X_n)^p]^{\frac{1}{p}} \EE[(\widetilde{X}_n)^p]^{\frac{p-1}{p}} \]
    giving the assertion (1) of the theorem (we use the fact that $\EE[(\widetilde{X}_n)^p] < \infty$). Then for the continuous case, we can use the same argument as in assertion (1) by considering a countable dense subset $D$ of $\RR_+$ such that $0 \in D$ and $t \in D$, and then applying the discrete case to the sequence $Y_n = X_{t^m_{n \wedge m}}$ for every fixed $m$. 

    \item The assertion (4) of the theorem follows from the assertion (3) applied to the submartingale $X_n = |Y_n|$. 
\end{itemize}
}

In what follows, you will find that most of the results stated in the continuous time case requiring the right-continuity of the sample paths of the process, which is not that hard to verify in most of the cases thanks to the following theorem, which gives a sufficient condition to find a modification with the desired regularity of the sample paths.

\thmpr{Sufficient condition for càdlàg modification}{cadlag_modif_existence}{
    Assume that the filtration $(\F_t)$ is right-continuous and complete. Let $X=(X_t)_{t \geq 0}$ be a supermartingale, such that the function  $t \mapsto \EE[X_t]$ is right-continuous. Then $X$ has a modification with càdlàg\footnote{refers to the french sentence "continue à droite et limite à gauche" which means "right-continuous and have left-limits"} sample paths, which is also an $(\F_t)$-supermartingale.
}{
    The proof of this theorem could be found in \parencite[thm.~3.18]{Stochastic}
}

\thmpr{Doob's Lp convergence theorem}{Doob_Lp_convergence}{
    Let $(X_n)_{n \in \NN}$ be a martingale and $p \in (1, \infty)$. Suppose that
    \[ \sup_{n \in \NN} \EE[|X_n|^p] < \infty, \]
    Then, $X_n$ converges a.s. and in $L^p$ to a random variable $X_\infty$ such that
    \[ \EE[|X_\infty|^p] = \sup_{n \in \NN} \EE[|X_n|^p] \]
    and we have
    \[ \EE[(\sup_{n \in \NN} |X_n|^p)] \leq \left( \frac{p}{p-1} \right)^p \EE[|X_\infty|^p]. \]
}{
    Since the martingale $(X_n)$ is bounded in $L^p$, it is also bounded in $L^1$, and we know from Theorem \ref{thm:as_conv_supermartingale} that $X_n$ converges a.s. to a limiting random variable $X_\infty$. Moreover, by setting $X_n^* = \sup_{1\leq k \leq n} |X_k|$, the theorem \ref{thm:doob_supermartingale_inequalities} shows that, for every $n \in \ZZ_+$,

    \begin{align*}
    \EE[(X_n^*)^p] \leq \left( \frac{p}{p-1} \right)^p \sup_{k \in \ZZ_+} \EE[|X_k|^p].
    \end{align*}

    Since $X_n^* \uparrow X_\infty^*$ when $n \uparrow \infty$, the monotone convergence theorem gives

    \begin{align*}
    \EE[(X_\infty^*)^p] \leq \left( \frac{p}{p-1} \right)^p \sup_{k \in \ZZ_+} \EE[|X_k|^p] < \infty
    \end{align*}

    and thus $X_\infty^* \in L^p$. Since all random variables $|X_n|^p$ are dominated by $|X_\infty^*|^p$, the dominated convergence theorem ensures that $X_n$ converges to $X_\infty$ in $L^p$. Finally, we know from proof of (3) in theorem \ref{thm:doob_supermartingale_inequalities} that the sequence $\EE[|X_n|^p]$ is increasing, and therefore

    \begin{align*}
    \EE[|X_\infty|^p] = \lim_{n \to \infty} \EE[|X_n|^p] = \sup_{n \in \ZZ_+} \EE[|X_n|^p].
    \end{align*}

}

\thmpr{Kolmogorov's criterion for a.s. convergence}{Kolmogorov_criterion}{
    Let $(X_n)_{n \in \NN}$ be a sequence of independent real random variables in $L^2$. Assume that $\EE[X_n] = 0$ for every $n \in \NN$. Then the following two conditions are equivalent:
    \begin{enumerate}
        \item $\displaystyle \sum_{n=0}^{\infty} \EE[(X_n)^2] < \infty$.
        \item  The series $\displaystyle \sum_{n=0}^{\infty} X_n$ converges a.s. and in $L^2$.
    \end{enumerate}
}{
    For every integer $n \geq 0$, set $\F_n = \sigma(X_0, X_1, \dots, X_n)$ and
    \[ M_n = \sum_{k=0}^n X_k. \]
    Since
    \[ \EE[M_{n+1} | \F_n] = \sum_{k=0}^n X_k + \EE[X_{n+1}] = \sum_{k=0}^n X_k, \]

    $(M_n)_{n \in \ZZ_+}$ is a martingale with respect to the filtration $(\F_n)_{n \in \ZZ_+}$. Furthermore, since $\EE[X_j X_k] = 0$ if $j \neq k$, one has

    \[\EE[(M_n)^2] = \sum_{k=0}^n \EE[(X_k)^2].\]

    Hence, if $M_n$ converges in $L^2$ as $n \to \infty$, condition (1) of the theorem must hold (we have just recovered a special case of a classical result of Hilbert space theory, see Riesz Fischer Theorem  \ref{thm:Riesz_Fischer} in the Appendix). Conversely, if (1) holds, the martingale $M_n$ is bounded in $L^2$, and by Theorem \ref{thm:Doob_Lp_convergence} we get that $M_n$ converges a.s. and in $L^2$. 
}

\subsection{Almost sure convergence}
\thmpr{Almost convergence of supermartingales}{as_conv_supermartingale}{
    Let $X$ be a supermartingale (with right-continuous sample paths if time is continuous). Assume that the collection $(X_n)_{n\in\NN}$ (resp. $(X_t)_{t \geq 0}$) is bounded in $L^1$. Then, it converge almost surely to a random variable $X_\infty\in L^1$.
}{
    We will prove it in the continuous time case, the discrete time case can be proven by similar arguments and it is left exercice for the reader. Let $D$ be a countable dense subset of $\RR_+$. Using lemma \ref{lem:Doob_upcrossing_inequality_continuous}, we have, for every $a < b$ and every $T\in D$, 
    \[ \EE\left[ M_{ab}^X(D \cap [0, T]) \right] \leq \frac{1}{b - a} \EE\left[ (X_T - a)^- \right] \]
    By Monotone convergence theorem, we get, for every $a<b$,
    \[ \EE\left[ M_{ab}^X(D) \right]  \leq \frac{1}{b - a} \sup_{t \geq 0} \EE\left[ (X_t - a)^- \right] < +\infty. \]
    Since the collection $(X_t)_{t \geq 0}$ is bounded in $L^1$, the right-hand side of the previous display is finite. Hence, a.s. for all rationals $a < b$, we have $M_{ab}^X(D) < \infty$. It follows from Lemma \ref{lem:limits_and_upcrossing_number_continuous} that the limit 
    \[ X_\infty := \lim_{D \ni t \to \infty} X_t \]
    exists a.s. in $[-\infty, +\infty]$. We can in fact exclude the values $-\infty$ and $+\infty$ since, by Fatou's lemma, 
    \[\EE[|X_\infty|] \leq \liminf_{D \ni t \to \infty} \EE[|X_t|] < +\infty\]
    and we get that $X_\infty \in L^1$. The right-continuity of $X$ samples (which we have not used yet) allows us to remove the restriction to $D$  in the limit.
}
\cor{Let $(X_n)_{n \in \NN}$ (resp. $(X_t)_{t \geq 0}$) be a nonnegative supermartingale. Then, $X_n$ (resp. $X_t$) converges a.s. as $n \to \infty$ (resp. $t \to \infty$). Its limit $X_\infty$ is in $L^1$ and we have $X_n \geq \EE[X_\infty \mid \F_n]$ for every $n \in \NN$. (resp. $X_t \geq \EE[X_\infty \mid \F_t]$ for every $t \geq 0$).
}

\paragraph{Backward supermartingales}
A useful variant of the previous theorem is the following one, which is stated in a backward form in discrete time setup, we will not prove it however proof could be found in \parencite[thm.~12.30]{Stochastic}. Since it will be used only in a proof step of supermartingale optional theorem in the continuous time case. 

A backward filtration is a sequence $(\F_n)_{n \in \ZZ_-}$ of sub-$\sigma$-fields of $\A$, which is indexed by the set $\ZZ_- = \{0, -1, -2, \dots\}$ of all nonpositive integers, and is such that $\F_n \subset \F_m$ whenever $n, m \in \ZZ_-$ and $n \leq m$. We then set

$$\F_{-\infty} = \bigcap_{n \in \ZZ_-} \F_n$$

which is also a sub-$\sigma$-field of $\A$. It is important to observe that, in contrast with the "forward case" studied in the previous sections, the $\sigma$-field $\F_n$ becomes smaller and smaller when $n \downarrow -\infty$ (the larger $|n|$ is, the smaller the $\sigma$-field $\F_n$).

In what follows, we fix a backward filtration $(\F_n)_{n \in \ZZ_-}$. Let $X = (X_n)_{n \in \ZZ_-}$ be a sequence of real random variables indexed by $\ZZ_-$. We say that $X$ is a backward martingale (resp. a backward supermartingale, a backward submartingale) if $X_n$ is $\F_n$-measurable and $\EE[|X_n|] < \infty$ for every $n \in \ZZ_-$, and if, for every $m, n \in \ZZ_-$ such that $n \leq m$,

$$X_n = \EE[X_m | \F_n] \quad (\text{resp. } X_n \geq \EE[X_m | \F_n], X_n \leq \EE[X_m | \F_n]).$$

\thmr{Backward supermartingale convergence}{backward_supermartingale_convergence}{
    Let $(X_n)_{n \in \ZZ_-}$ be a backward supermartingale. Suppose that
    $$\sup_{n \in \ZZ_-} \EE[|X_n|] < \infty. $$
    Then the sequence $(X_n)_{n \in \ZZ_-}$ is uniformly integrable and converges a.s. and in $L^1$ to a random variable $X_\infty$ as $n \to -\infty$. Moreover for every $n \in \ZZ_-$,
    $$\EE[X_n | \F_{-\infty}] \leq X_\infty,$$
    and equality holds in the last display if $(X_n)_{n \in \ZZ_-}$ is a backward martingale.
}


\subsection{Martingale Closure} \label{sec:uniform_integrability}
\defnr{Closed Martingale}{closed_martingale}{A martingale $(X_n)_{n\in\NN}$ (resp. $(X_t)_{t \geq 0}$ ) is said to be \textit{closed} if there exists a random variable $Z \in L^1$ such that, for every $n\in\NN$ (resp. $t \geq 0$),
\[ \begin{cases}
    X_n = \EE[Z \mid \F_n] & \text{For the discrete time} \\ 
    X_t = \EE[Z \mid \F_t] & \text{For the continuous time} 
\end{cases} \]

}
\defnr{Uniform integrability}{uniform_integrability}{A collection $(X_i)_{i \in I}$ of random variables in $L^1(\Omega, \A, \PP)$ is said to be \textit{uniformly integrable} (u.i. in short) if
\[ \lim_{a \to +\infty} \left( \sup_{i \in I} \EE[|X_i| \1_{\{|X_i| > a\}}] \right) = 0. \]
}
A uniformly integrable collection $(X_i)_{i \in I}$ is bounded in $L^1$. To verify this, take $a$ large enough so that
\[ \sup_{i \in I} \EE[|X_i| \1_{\{|X_i| > a\}}] \leq 1 \]
and write $\EE[|X_i|] \leq \EE[|X_i| \1_{\{|X_i| \leq a\}}] + \EE[|X_i| \1_{\{|X_i| > a\}}] \leq a + 1$. The converse is false: a collection $(X_i)_{i \in I}$ which is bounded in $L^1$ may not be uniformly integrable.


\proppr{characterstic_UI}{Let $(X_i)_{i \in I}$ be a bounded subset of $L^1$. The following two conditions are equivalent:
\begin{enumerate}
    \item The collection $(X_i)_{i \in I}$ is u.i.
    \item For every $\varepsilon > 0$, one can find $\delta > 0$ such that, for every event $A \in \A$ of probability $\PP(A) < \delta$, one has
    \[ \forall i \in I, \quad \EE[|X_i|\1_A] < \varepsilon. \]
\end{enumerate}
}{
    (1)$\Rightarrow$(2) Let $\varepsilon > 0$. We can first fix $a > 0$ such that
    \[  \sup_{i \in I} \EE[|X_i| \1_{\{|X_i|>a\}}] < \frac{\varepsilon}{2}.     \]
    If we set $\delta = \varepsilon / (2a)$, then the condition $\PP(A) < \delta$ implies, that, for every $i \in I$,
    \[
    \EE[|X_i| \1_A] \leq \EE[|X_i| \1_{A \cap \{|X_i| \leq a\}}] + \EE[|X_i| \1_{\{|X_i| > a\}}] \leq a \PP(A) + \frac{\varepsilon}{2} < \varepsilon.
    \]
    \noindent (2)$\Rightarrow$(1) Set $C = \sup_{i \in I} \EE[|X_i|]$. By the Markov inequality, we have for every $a > 0$,
    \[ \forall i \in I, \quad \PP(|X_i| > a) \leq \frac{C}{a}. \]
    Let $\varepsilon > 0$ and choose $\delta > 0$ so that the property stated in (2) holds. Then, if $a$ is large enough so that $C/a < \delta$, we have $\PP(|X_i| > a) < \delta$ for every $i \in I$, and thus, by (2),
    \[ \forall i \in I, \quad \EE[|X_i| \1_{\{|X_i| > a\}}] < \varepsilon. \]
    We conclude that the collection $(X_i)_{i \in I}$ is u.i.}


\thmpr{Vitali convergence theorem}{Vitali_convergence_theorem}{
    Let $(X_n)_{n \in \NN}$ be a sequence of random variables in $L^1$ that converges to $X_\infty$ in probability. The following two conditions are equivalent:
    \begin{enumerate}
        \item The sequence $(X_n)_{n \in \NN}$ converges to $X_\infty$ in $L^1$.
        \item The sequence $(X_n)_{n \in \NN}$ is uniformly integrable.
    \end{enumerate}
}{
\begin{itemize}
    \item (1)$\Rightarrow$(2) If the sequence $(X_n)_{n \in \NN}$ converges in $L^1$, then it is bounded in $L^1$. Then, let $\varepsilon > 0$. We can choose $N$ large enough so that, for every $n \geq N$,
    \[ \EE[|X_n - X_N|] < \frac{\varepsilon}{2}. \]
    Since the finite set $\{X_1, \dots, X_N\}$ is u.i., Proposition \ref{prop:characterstic_UI} allows us to choose $\delta > 0$ such that, for every event $A$ of probability $\PP(A) < \delta$, we have
    \[\forall n \in \{1, \dots, N\}, \quad \EE[|X_n|\1_A] < \frac{\varepsilon}{2}.\]
    Then, if $n > N$, we have also
    \[\EE[|X_n|\1_A] \leq \EE[|X_N|\1_A] + \EE[|X_n - X_N|] < \varepsilon.\]
    By Proposition \ref{prop:characterstic_UI}, the sequence $(X_n)_{n \in \NN}$ is uniformly integrable.
    \item (2)$\Rightarrow$(1) Using the characterization of uniform integrability in Proposition \ref{prop:characterstic_UI}(ii), it is straightforward to verify that the collection $(X_n - X_m)_{n,m \in \NN}$ is also u.i. Hence, if $\varepsilon > 0$ is fixed, we can choose $a > 0$ large enough so that, for every $m, n \in \NN$,
    \[ \EE[|X_n - X_m|\1_{\{|X_n - X_m| > a\}}] < \varepsilon. \]
    Then, for every $m, n \in \NN$,
    \begin{align*}
        \EE[|X_n - X_m|] &\leq \EE[|X_n - X_m|\1_{\{|X_n - X_m| \leq \varepsilon\}}] + \EE[|X_n - X_m|\1_{\{\varepsilon < |X_n - X_m| \leq a\}}] \\
        &\quad + \EE[|X_n - X_m|\1_{\{|X_n - X_m| > a\}}] \\
        &\leq 2\varepsilon + a \, \PP(|X_n - X_m| > \varepsilon).
    \end{align*}
    The convergence in probability of the sequence $(X_n)_{n \in \NN}$ implies that
    \[\PP(|X_n - X_m| > \varepsilon) \leq \PP\left(|X_n - X_\infty| > \frac{\varepsilon}{2}\right) + \PP\left(|X_m - X_\infty| > \frac{\varepsilon}{2}\right) \xrightarrow[n,m \to \infty]{} 0. \]
    It follows that
    \[ \limsup_{m,n \to \infty} \EE[|X_n - X_m|] \leq 2\varepsilon \]
    and, since $\varepsilon$ was arbitrary, this shows that the sequence $(X_n)_{n \in \NN}$ is Cauchy in $L^1$, hence converges in $L^1$.
\end{itemize}

}
\thmpr{Doob's Uniform Integrability convergence theorem}{Doob_UI_convergence}{
Let $X$ be a martingale (with right-continuous sample paths if time is continuous). Then the following properties are equivalent:
\begin{enumerate}
    \item $X$ is closed;
    \item the collection $(X_n)_{n\in\NN}$ ($(X_t)_{t \geq 0}$ if time is continuous) is uniformly integrable;
    \item $X_n$ (resp. $X_t$) converges a.s. and in $L^1$ as $n \to \infty$ (resp. $t \to \infty$).
\end{enumerate}
Moreover, if these properties hold, we have  $X_n = \EE[X_\infty \mid \F_n]$ (resp. $X_t = \EE[X_\infty \mid \F_t]$) for every $n \in \NN$ (resp. $t \geq 0$), where $X_\infty \in L^1$ is the a.s. limit of $X_n$ (resp. $X_t$) as $n \to \infty$ (resp. $t \to \infty$).
}{
    \begin{itemize}
        \item (1)$\Rightarrow$(2): if $Z\in L^1$, the collection of all random variable $\EE[Z|\G]$ when $\G$ varies among all sub-$\sigma$-fields of $\F_\infty$ is u.i. 
        \item (2)$\Rightarrow$(3): If (2) holds, then $(X_n)_{n\in\NN}$ (resp. $(X_t)_{t \geq 0}$) is bounded in $L^1$, and thus converges a.s. to some random variable $X_\infty$ by Theorem \ref{thm:as_conv_supermartingale}.  By uniform integrability, the latter convergence also holds in $L^1$.
        \item (3)$\Rightarrow$(1): If (3) holds, for every $m \in \NN$ (resp. $s \geq 0$), we can pass to the limit $n\to +\infty$ (resp. $t\to +\infty$) in the equality, , $X_m = \EE(X_n \mid \F_m)$ (resp. $X_s = \EE(X_t \mid \F_s)$) and using the theorem \ref{prop:conditional_expectation_limits_properties} we get $X_m = \EE(X_\infty \mid \F_m)$ (resp. $X_s = \EE(X_\infty \mid \F_s)$), which gives the desired result.
    \end{itemize}
}

\cor{Let $Z \in L^1(\Omega, \A, \PP)$. The martingale $X_n = \EE[Z \mid \F_n]$ (resp. $X_t = \EE[Z \mid \F_t]$) converges a.s. and in $L^1$ to $X_\infty = \EE[Z \mid \F_\infty]$.}


\section{Optional stopping theorems}


\subsection{The theorem for martingales}
\thmpr{Optional stopping theorem -- martingale}{optional_stopping_martingale}{
    Let $(X_n)_{n\in\NN}$ (resp. $(X_t)_{t \geq 0}$) be a uniformly integrable martingale (with right-continuous sample paths if time is continnuous). 
    Let $S$ and $T$ be two stopping times with $S \leq T$. Then $X_S$ and $X_T$ are in $L^1$ and
    \[     X_S = \EE[X_T \mid \F_S]. \]
}{
    We begins with the discrete time case. From the preceding section, we know that the martingale $(X_n)_{n \in \NN}$ is closed and $X_n = \EE[X_\infty | \F_n]$ for every $n$. Let us verify that $X_T \in L^1$:

    \begin{align*}
    \EE[|X_T|] 
        &= \sum_{n=0}^{\infty} \EE[\1_{\{T=n\}} |X_n|] + \EE[\1_{\{T=\infty\}} |X_\infty|] \\
        &= \sum_{n=0}^{\infty} \EE[\1_{\{T=n\}} |\EE[X_\infty | \F_n]|] + \EE[\1_{\{T=\infty\}} |X_\infty|] \\
        &\le \sum_{n=0}^{\infty} \EE[\1_{\{T=n\}} \EE[|X_\infty| | \F_n]] + \EE[\1_{\{T=\infty\}} |X_\infty|] \\
        &= \sum_{n=0}^{\infty} \EE[\1_{\{T=n\}} |X_\infty|] + \EE[\1_{\{T=\infty\}} |X_\infty|] \\
        &= \EE[|X_\infty|] < \infty.
    \end{align*}

    Then, if $A \in \F_T$,

    \begin{align*}
        \EE[\1_A X_T] &= \sum_{n \in \NN \cup \{\infty\}} \EE[\1_{A \cap \{T=n\}} X_T] \\
        &= \sum_{n \in \NN \cup \{\infty\}} \EE[\1_{A \cap \{T=n\}} X_n] \\
        &= \sum_{n \in \NN \cup \{\infty\}} \EE[\1_{A \cap \{T=n\}} X_\infty] \\
        &= \EE[\1_A X_\infty].
    \end{align*}

    In the first equality, we use the fact that $X_T \in L^1$ to justify the interchange of sum and expectation via the Fubini theorem. In the third equality, we use the properties $X_n = \EE[X_\infty | \F_n]$ and $A \cap \{T = n\} \in \F_n$ for $A \in \F_T$. Since $X_T$ is $\F_T$-measurable and in $L^1$, the preceding display shows that $X_T$ satisfies the characteristic property of the conditional expectation $\EE[X_\infty | \F_T]$.

    Finally, if $S$ and $T$ are two stopping times such that $S \le T$, the property $\F_S \subset \F_T$ (Proposition \ref{prop:two_stopping_times}) and Theorem of nested \sfields in conditional expectation \ref{thm:conditional_nested_sigma_fields} imply that
    \[     X_S = \EE[X_\infty | \F_S] = \EE[\EE[X_\infty | \F_T] | \F_S] = \EE[X_T | \F_S].     \]

    For the continuous time case, Set, for every integer $n \geq 0$,
    \[   T_n=  \frac{\lceil 2^n T \rceil}{2^n}\cdot \1_{\{T<\infty\}}+\infty \cdot \1_{\{T=\infty\}}, \quad n=0,1,2, \ldots     \]
    and similarly
    \[    S_n=  \frac{\lceil 2^n S \rceil}{2^n}\cdot \1_{\{S<\infty\}}+\infty \cdot \1_{\{S=\infty\}}, \quad n=0,1,2, \ldots \]
    By Proposition \ref{prop:majorant_is_stoping_time}, $(T_n)$ and $(S_n)$ are two sequences of stopping times that decrease respectively to $T$ and to $S$. Moreover, we have $S_n \leq T_n$ for every $n \geq 0$.

    Now observe that, for every fixed $n$, $2^n S_n$ and $2^n T_n$ are stopping times of the discrete filtration $\H_k^{(n)} := \F_{k/2^n}$, and $Y_k^{(n)} := X_{k/2^n}$ is a discrete martingale with respect to this filtration. From the discrete case we proven before, we get that $Y_{2^n S_n}^{(n)}$ and $Y_{2^n T_n}^{(n)}$ are in $L^1$, and
    \[ X_{S_n} = Y_{2^n S_n}^{(n)} = \EE[Y_{2^n T_n}^{(n)} \mid \H_{2^n S_n}^{(n)}] = \EE[X_{T_n} \mid \F_{S_n}] \]
    (here we need to verify that $\H_{2^n S_n}^{(n)} = \F_{S_n}$, but this is straightforward).

    Let $A \in \F_S$. Since $\F_S \subset \F_{S_n}$, we have $A \in \F_{S_n}$ and thus
    \[ \EE[\1_A X_{S_n}] = \EE[\1_A X_{T_n}].\]
    By the right-continuity of sample paths, we get a.s.
    \[ X_S = \lim_{n \to \infty} X_{S_n}, \quad X_T = \lim_{n \to \infty} X_{T_n}. \]
    These limits also hold in $L^1$. Indeed, thanks again to the optional stopping theorem for uniformly integrable discrete martingales, we have $X_{S_n} = \EE[X_\infty \mid \F_{S_n}]$ for every $n$, and thus the sequence $(X_{S_n})$ is uniformly integrable (and the same holds for the sequence $(X_{T_n})$).

    The $L^1$-convergence implies that $X_S$ and $X_T$ belong to $L^1$, and also allows us to pass to the limit $n \to \infty$ in the equality $\EE[\1_A X_{S_n}] = \EE[\1_A X_{T_n}]$ in order to get
    \[
    \EE[\1_A X_S] = \EE[\1_A X_T].
    \]
    Since this holds for every $A \in \F_S$, and since the variable $X_S$ is $\F_S$-measurable (by the remarks before the theorem), we conclude that
    \[
    X_S = \EE[X_T \mid \F_S],
    \]
    which completes the proof.
}

In particular, for every stopping time $S$, we have
\[     X_S = \EE[X_\infty \mid \F_S],\]
and
\[\EE[X_S] = \EE[X_\infty] = \EE[X_0]. \]

\corp{
    If $T$ and $S$ are stopping time such that $S \leq T$, and both $S$ and $T$ are bounded, then $X_S$ and $X_T$ are in $L^1$ and
    \[
        X_S = \EE[X_T \mid \F_S].
    \]
}{
Let $a \geq 0$ such that $S \leq T \leq a$. We apply Optional Stopping Theorem \ref{thm:optional_stopping_martingale} to the martingale $\{X_{n \wedge a}\}_{n \geq 0}$ (resp. $\{X_{t \wedge a}\}_{t \geq 0}$) which is closed by $X_a$.
}

\corp{ Let $(X_n)_{n\in\NN}$ (resp. $(X_t)_{t \geq 0}$) be a martingale (with right-continuous sample paths if time is continuous), and let $T$ be a stopping time.
\begin{enumerate}
    \item The process $(X_{n \wedge T})_{n\in\NN}$  (resp. $(X_{t \wedge T})_{t \geq 0}$) is still a martingale.
    \item Suppose in addition that the martingale $(X_n)_{n\in\NN}$ (resp. $(X_t)_{t \geq 0}$) is uniformly integrable. Then the process $(X_{n \wedge T})_{n\in\NN}$ (resp. $(X_{t \wedge T})_{t \geq 0}$) is also a uniformly integrable martingale, and more precisely we have,
    \[
    \begin{cases}
        X_{n \wedge T} = \EE[X_T \mid \F_{n \wedge T}], \quad n \in \NN
        & \text{for the discrete case}, \\[4pt]
        X_{t \wedge T} = \EE[X_T \mid \F_{t \wedge T}], \quad t \ge 0
        & \text{for the continuous case}.
    \end{cases}
    \]
\end{enumerate}
}{
    We will prove it for the continuous time case, the discrete time case can be proven by similar arguments and it is left exercice for the reader. 
    We start with the proof of (2). Note that $t \wedge T$ is a stopping time by property (f) of stopping times. By Theorem \ref{thm:optional_stopping_martingale}, $X_{t \wedge T}$ and $X_T$ are in $L^1$, and we also know that $X_{t \wedge T}$ is $\F_{t \wedge T}$-measurable, hence $\F_t$-measurable since $\F_{t \wedge T} \subset \F_t$. So in order to get formula in (2), it is enough to prove that, for every $A \in \F_t$,

    $$\EE[\1_A X_T] = \EE[\1_A X_{t \wedge T}].$$

    Let us fix $A \in \F_t$. First, we have trivially

    \[ \EE[\1_{A \cap \{T \leq t\}} X_T] = \EE[\1_{A \cap \{T \leq t\}} X_{t \wedge T}]. \tag{*} \]

    On the other hand, by Theorem \ref{thm:optional_stopping_martingale}, we have

    $$X_{t \wedge T} = \EE[X_T \mid \F_{t \wedge T}],$$

    and we notice that we have both $A \cap \{T > t\} \in \F_t$ and $A \cap \{T > t\} \in \F_T$ (the latter as a straightforward consequence of the definition of $\F_T$), so that $A \cap \{T > t\} \in \F_t \cap \F_T = \F_{t \wedge T}$. Using the preceding display, we obtain

    $$\EE[\1_{A \cap \{T > t\}} X_T] = \EE[\1_{A \cap \{T > t\}} X_{t \wedge T}].$$

    By adding this equality to (*), we get the desired result.

    To prove (1), we just need to apply (2) to the (uniformly integrable) martingale $(X_{t \wedge a})_{a \geq 0}$, for any choice of $a \geq 0$. 
}


\subsection{The theorem for supermartingales}

\thmpr{Optional stopping theorem - supermartingale}{optional_stopping_supermartingale}{
Let $(Z_n)_{n\in\NN}$ (resp. $(Z_t)_{t \geq 0}$) be a nonnegative supermartingale (with right-continuous sample paths if time is continnuous). Let $U$ and $V$ be two stopping times such that $U \le V$. Then $Z_U$ and $Z_V$ are in $L^1$, and
    \[ Z_U \ge \EE[Z_V \mid \F_U]. \]
}{
    Let's begin with the discrete time case. We first note that $(X_n)_{n \in \NN}$ is bounded in $L^1$. Theorem \ref{thm:as_conv_supermartingale} then shows that $X_n$ converges a.s. to $X_\infty$, so that the definition of $X_T$ makes sense for any stopping time, even on the event $\{T = \infty\}$. 
    If $T$ is a bounded stopping time, we know that $\EE[X_T] \leq \EE[X_0]$ (Lemma \ref{lem:optional_stopping_theorem_v1}). Fatou's lemma then shows that, for any stopping time $T$, 
    $$\EE[X_T] \leq \liminf_{k \to \infty} \EE[X_{T \wedge k}] \leq \EE[X_0]$$
    and thus $X_T \in L^1$.

    Then let $S$ and $T$ be two stopping times such that $S \leq T$. Let us first assume that $T \leq N$ for some fixed integer $N$. According to Lemma \ref{lem:increasing_stopping_time_submartingale}, we have then $\EE[X_S] \geq \EE[X_T]$. More generally, let $A \in \F_S$, and set

    $$S^A(\omega) = 
    \begin{cases} 
        S(\omega) & \text{if } \omega \in A, \\ 
        N & \text{if } \omega \notin A. 
    \end{cases}$$

    It is straightforward to verify that $S^A$ is a stopping time (just note that $\{S^A \leq n\} = \Omega$ if $n \geq N$, and $\{S^A \leq n\} = A \cap \{S \leq n\} \in \F_n$ if $n < N$, using the definition of $\F_S$). We define $T^A$ in a similar manner, and note that $T^A$ is also a stopping time since $A \in \F_S \subset \F_T$. Since $S^A \leq T^A$, we have $\EE[X_{S^A}] \geq \EE[X_{T^A}]$, and it follows that

    $$\EE[X_S \1_A] \geq \EE[X_T \1_A].$$

    Let us come back to the general case where $S$ and $T$ satisfy $S \leq T$ but are no longer assumed to be bounded. Let $A \in \F_S$, and $k \in \NN$. Then $A \cap \{S \leq k\}$ belongs both to $\F_k$ (by the definition of $\F_S$) and to $\F_S$ (recall that $S$ is $\F_S$-measurable). From Proposition \ref{prop:two_stopping_times}, $A \cap \{S \leq k\}$ belongs to $\F_{S \wedge k}$. Hence, by applying the bounded case to the stopping times $S \wedge k$ and $T \wedge k$, we get

    $$\EE[X_S \1_{A \cap \{S \leq k\}}] = \EE[X_{S \wedge k} \1_{A \cap \{S \leq k\}}] \geq \EE[X_{T \wedge k} \1_{A \cap \{S \leq k\}}].$$

    By monotone convergence,

    $$\lim_{k \to \infty} \uparrow \EE[X_S \1_{A \cap \{S \leq k\}}] = \EE[X_S \1_{A \cap \{S < \infty\}}].$$

    On the other hand, Fatou's lemma gives

    $$\liminf_{k \to \infty} \EE[X_{T \wedge k} \1_{A \cap \{S \leq k\}}] \geq \EE[X_T \1_{A \cap \{S < \infty\}}].$$

    So we have obtained that

    $$\EE[X_S \1_{A \cap \{S < \infty\}}] \geq \EE[X_T \1_{A \cap \{S < \infty\}}].$$

    Since $X_S = X_T = X_\infty$ on the event $\{S = \infty\}$, it follows that

    $$\EE[X_S \1_A] \geq \EE[X_T \1_A] = \EE[\EE[X_T \mid \F_S] \1_A].$$

    Since this holds for every $A \in \F_S$, and both $X_S$ and $\EE[X_T \mid \F_S]$ are $\F_S$-measurable, it follows that $X_S \geq \EE[X_T \mid \F_S]$ a.s. (apply the last display with the set $A = \{X_S < \EE[X_T \mid \F_S]\}$). This completes the proof of the theorem in the discrete case.

    Now the continuous time case, Firstly, we make the extra assumption that $U$ and $V$ are bounded and we verify that we then have $\EE[Z_U] \ge \EE[Z_V]$. Let $p \ge 1$ be an integer such that $U \le p$ and $V \le p$. For every integer $n \ge 0$, set
    \[
    U_n = \sum_{k=0}^{p2^n-1} \frac{k+1}{2^n} \1_{\{k2^{-n} < U \le (k+1)2^{-n}\}}, \quad V_n = \sum_{k=0}^{p2^n-1} \frac{k+1}{2^n} \1_{\{k2^{-n} < V \le (k+1)2^{-n}\}}
    \]
    in such a way (by Proposition \ref{prop:majorant_is_stoping_time}) that $(U_n)$ and $(V_n)$ are two sequences of bounded stopping times that decrease respectively to $U$ and $V$, and additionally we have $U_n \le V_n$ for every $n \ge 0$. The right-continuity of sample paths ensures that $Z_{U_n} \to Z_U$ and $Z_{V_n} \to Z_V$ a.s. as $n \to \infty$. Then, by the optional stopping theorem for discrete supermartingales in the case of bounded stopping times (which we just proved), with respect to the filtration $(\F_{k/2^{n+1}})_{k\ge0}$, we have for every $n \ge 0$,
    \[ Z_{U_{n+1}} \ge \EE[Z_{U_n} \mid \F_{U_{n+1}}]. \]
    Setting $Y_n = Z_{U_{-n}}$ and $\H_n = \F_{U_{-n}}$, for every integer $n \le 0$, we get that the sequence $(Y_n)_{n\le0}$ is a backward supermartingale with respect to the filtration $(\H_n)_{n\le0}$. Since, for every $n \ge 0$, $\EE[Z_{U_n}] \le \EE[Z_0]$ (by another application of the discrete optional stopping theorem), the sequence $(Y_n)_{n\le0}$ is bounded in $L^1$, and by the convergence theorem for backward supermartingales in Theorem \ref{thm:backward_supermartingale_convergence}, it converges in $L^1$. Hence the convergence of $Z_{U_n}$ to $Z_U$ also holds in $L^1$ and similarly the convergence of $Z_{V_n}$ to $Z_V$ holds in $L^1$. Since $U_n \le V_n$, by yet another application of the discrete optional stopping theorem, we have $\EE[Z_{U_n}] \ge \EE[Z_{V_n}]$. Using the $L^1$-convergence of $Z_{U_n}$ and $Z_{V_n}$ we can pass to the limit $n \to \infty$ and obtain that $\EE[Z_U] \ge \EE[Z_V]$ as claimed.

    Let us prove the statement of the theorem (no longer assuming that $U$ and $V$ are bounded). By the first step of the proof applied to the stopping times $0$ and $U \wedge p$, we have $\EE[Z_{U \wedge p}] \le \EE[Z_0]$ for every $p \ge 1$, and Fatou's lemma gives $\EE[Z_U] \le \EE[Z_0] < \infty$ and similarly $\EE[Z_V] < \infty$. Fix $A \in \F_U \subset \F_V$ and recall our notation $U^A$ for the stopping time defined by $U^A(\omega) = U(\omega)$ if $\omega \in A$ and $U^A(\omega) = \infty$ otherwise (cf. property (d) of stopping times). By the first part of the proof, we have, for every $p \ge 1$,
    \[ \EE[Z_{U^A \wedge p}] \ge \EE[Z_{V^A \wedge p}]. \]
    By writing each of these two expectations as a sum of expectations over the sets $A^c$, $A \cap \{U \le p\}$ and $A \cap \{U > p\}$, and noting that $U > p$ implies $V > p$, we get
    \[ \EE[Z_U \1_{A \cap \{U \le p\}}] \ge \EE[Z_V \1_{A \cap \{U \le p\}}]. \]
    Letting $p \to \infty$ then gives
    \[ \EE[Z_U \1_{A \cap \{U < \infty\}}] \ge \EE[Z_V \1_{A \cap \{U < \infty\}}]. \]

    is trivial and by adding it to the preceding display, we obtain
    \[ \EE[Z_U \1_A] \ge \EE[Z_V \1_A] = \EE[\EE[Z_V \mid \F_U] \1_A]. \]
    Since this holds for every $A \in \F_U$ and $Z_U$ is $\F_U$-measurable, the desired result follows.
}

\textbf{Remark} \quad
This implies that $\EE[Z_U] \ge \EE[Z_V]$, and since $Z_U = Z_V = Z_\infty$ on the event $\{U = \infty\}$, it also follows that
\[
\EE[\1_{\{U < \infty\}} Z_U]
\ge
\EE[\1_{\{U < \infty\}} Z_V]
\ge
\EE[\1_{\{V < \infty\}} Z_V].
\]



\section{Finite Variation Processes} \label{sec:finite-variation-processes}
\subsection{Finite Variation Functions}
In our discussion of functions with finite variation, we restrict our attention to \textit{continuous} functions, as this is the case of interest in the subsequent developments. Recall that a \textbf{signed measure} on a compact interval $[0, T]$ is the difference of two finite positive measures on $[0, T]$.
\defnr{Finite Variation function }{finit_var_func}{
    Let $T \ge 0$. A continuous function $a : [0, T] \longrightarrow \RR$ such that $a(0) = 0$ is said to have \textit{finite variation} if there exists a signed measure $\mu$ on $[0, T]$ such that $a(t) = \mu([0, t])$ for every $t \in [0, T]$.
}
The measure $\mu$ is then determined uniquely by $a$. Since $a$ is continuous and $a(0) = 0$, it follows that $\mu$ has no atoms\footnote{A measure $\mu$ has an atom $A$ if $\mu(A) > 0$ and for every $B \in \A$ with $B \subseteq A$, either $\mu(B) = 0$ or $\mu(B) = \mu(A)$.}.

By the Jordan Decomposition Theorem, any signed measure can be split uniquely as $\mu = \mu^+ - \mu^-$, where $\mu^+$ and $\mu^-$ are mutually singular positive measures, ensuring that the structural decomposition of such restricted systems is completely rigid and unique.
The decomposition of as a difference of two finite positive measures on $[0, T]$
is not unique, but there exists a unique decomposition $\mu = \mu^+ - \mu^-$ such that  $\mu^+$ and $\mu^-$ are supported on disjoint Borel sets ($D_+$ and $D_-$).
We write $|\mu|$ for the (finite) positive measure $|\mu| = \mu_+ + \mu_-$. The measure $|\mu|$ is called the \textit{total variation} of $a$. We have $|\mu(A)| \leq |\mu|(A)$ for every $A \in \B([0, T])$. Moreover, the Radon-Nikodym derivative of $\mu$ with respect to $|\mu|$ is
\[ \frac{d\mu}{d|\mu|} = \1_{D_+} - \1_{D_-}.\]
The fact that $a(t) = \mu_+([0, t]) - \mu_-([0, t])$ shows that $a$ is the difference of two monotone nondecreasing continuous functions that vanish at 0 (since $\mu$ has no atoms, the same holds for $\mu_+$ of $\mu_-$). Conversely, the difference of two monotone nondecreasing continuous functions that vanish at $0$ has finite variation in the sense of the previous definition. Indeed, this follows from the well-known fact that the formula $g(t) = \theta([0, t]), t \in [0, T]$ induces a bijection between monotone nondecreasing right-continuous functions $g : [0, T] \longrightarrow \RR_+$ and finite positive measures $\theta$ on $[0, T]$.

Let $f : [0, T] \longrightarrow \RR$ be a measurable function such that $\int_{[0,T]} |f(s)| \, |\mu|(ds) < \infty$. We set
\[
\begin{aligned}
\int_0^T f(s) \, da(s) &= \int_{[0,T]} f(s) \, \mu(ds), \\
\int_0^T f(s) \, |da(s)| &= \int_{[0,T]} f(s) \, |\mu|(ds).
\end{aligned}
\]

Then the bound
\[
\left| \int_0^T f(s) \, da(s) \right| \leq \int_0^T |f(s)| \, |da(s)|
\]
holds. By restricting $a$ to $[0, t]$ (which amounts to restricting $\mu$, $\mu_+$, $\mu_-$), we can define $\int_0^t f(s) \, da(s)$ for every $t \in [0, T]$, and we observe that the function $t \mapsto \int_0^t f(s) \, da(s)$ also has finite variation on $[0, T]$ (the associated measure is just $\mu'(ds) = f(s)\mu(ds)$).

\proppr{characterization_variation_func}{
For every $t \in (0, T]$,
\[\int_0^t |da(s)| = \sup \left\{ \sum_{i=1}^p |a(t_i) - a(t_{i-1})| \right\},\] 
where the supremum is over all subdivisions $0 = t_0 < t_1 < \dots < t_p = t$ of $[0, t]$. More precisely, for any increasing sequence $0 = t_0^n < t_1^n < \dots < t_{p_n}^n = t$ of subdivisions of $[0, t]$ whose mesh tends to 0, we have
\[ \lim_{n \to \infty} \sum_{i=1}^{p_n} |a(t_i^n) - a(t_{i-1}^n)| = \int_0^t |da(s)|. \]
}{
Clearly, it is enough to treat the case $t = T$. The inequality $\geq$ in the first assertion is very easy since, for any subdivision $0 = t_0 < t_1 < \dots < t_p = T$ of $[0, T]$,
\[ |a(t_i) - a(t_{i-1})| = |\mu((t_{i-1}, t_i])| \leq |\mu|((t_{i-1}, t_i]), \quad \forall i \in \{1, \dots, p\},\]
and
\[\sum_{i=1}^p |\mu|((t_{i-1}, t_i]) = |\mu|([0, T]) = \int_0^T |da(s)|. \] 

In order to get the reverse inequality, it suffices to prove the second assertion. So we consider an increasing sequence $0 = t_0^n < t_1^n < \dots < t_{p_n}^n = T$ of subdivisions of $[0, T]$, whose mesh $\max\{t_i^n - t_{i-1}^n : 1 \leq i \leq p_n\}$ tends to 0. Although we are proving a "deterministic" result, we will use a martingale argument. Leaving aside the trivial case where $|\mu| = 0$, we introduce the probability space $\Omega = [0, T]$, which is equipped with the Borel $\sigma$-field $\B = \B([0, T])$ and the probability measure $P(ds) = (|\mu|([0, T]))^{-1} |\mu|(ds)$. On this probability space, we consider the discrete filtration $(\B_n)_{n \geq 0}$ such that, for every integer $n \geq 0$, $\B_n$ is the $\sigma$-field generated by the intervals $(t_{i-1}^n, t_i^n]$, $1 \leq i \leq p_n$. We then set
\[ X(s) = \1_{D_+}(s) - \1_{D_-}(s) = \frac{d\mu}{d|\mu|}(s),\]
and, for every $n \geq 0$,
\[ X_n = \EE[X \mid \B_n]. \]

Properties of conditional expectation show that $X_n$ is constant on every interval $(t_{i-1}^n, t_i^n]$ and takes the value
\[ \frac{\mu((t_{i-1}^n, t_i^n])}{|\mu|((t_{i-1}^n, t_i^n])} = \frac{a(t_i^n) - a(t_{i-1}^n)}{|\mu|((t_{i-1}^n, t_i^n])} \]
on this interval. On the other hand, the sequence $(X_n)$ is a closed martingale, with respect to the discrete filtration $(\B_n)$. Since $X$ is measurable with respect to $\B = \bigvee_n \B_n$, this martingale converges to $X$ in $L^1$, by the convergence theorem for closed discrete martingales \ref{thm:Doob_UI_convergence}. In particular,
\[ \lim_{n \to \infty} \EE[|X_n|] = \EE[|X|] = 1, \]
where the last equality is clear since $|X(s)| = 1$, $|\mu|(ds)$ a.e. The desired result follows by noting that
\[ \EE[|X_n|] = (|\mu|([0, T]))^{-1} \sum_{i=1}^{p_n} |a(t_i^n) - a(t_{i-1}^n)|, \]
and recalling that $|\mu|([0, T]) = \int_0^T |da(s)|$.
}
\textbf{Remark} In the usual presentation of functions with finite variation, one starts from the property that the supremum in the first display of the proposition is finite.

We now give a useful approximation lemma for the integral of a continuous
function with respect to a function with finite variation

\lempr{Approximation for integral w.r.t finite variation function}{finite_variation_integral_proxy}{
If $f : [0, T] \longrightarrow \RR$ is a continuous function, and if $0 = t_0^n < t_1^n < \dots < t_{p_n}^n = T$ is a sequence of subdivisions of $[0, T]$ whose mesh tends to $0$, we have
\[ \int_0^T f(s) \, da(s) = \lim_{n \to \infty} \sum_{i=1}^{p_n} f(t_{i-1}^n) \left(a(t_i^n) - a(t_{i-1}^n)\right). \]
}{
Let $f_n$ be defined on $[0, T]$ by $f_n(s) = f(t_{i-1}^n)$ if $s \in (t_{i-1}^n, t_i^n]$, $1 \leq i \leq p_n$, and $f_n(0) = f(0)$. Then,
\[ \sum_{i=1}^{p_n} f(t_{i-1}^n) \left(a(t_i^n) - a(t_{i-1}^n)\right) = \int_{[0,T]} f_n(s) \, \mu(ds), \]
and the desired result follows by dominated convergence since $f_n(s) \longrightarrow f(s)$ as $n \to \infty$, for every $s \in [0, T]$. 
}

We say that a function $a : \RR_+ \longrightarrow \RR$ is a \textit{finite variation function on $\RR_+$} if the restriction of $a$ to $[0, T]$ has finite variation on $[0, T]$, for every $T > 0$. Then there is a unique $\sigma$-finite (positive) measure on $\RR_+$ whose restriction to every interval $[0, T]$ is the total variation measure of the restriction of $a$ to $[0, T]$, and we write
$$
\int_0^\infty f(s) \, |da(s)|
$$
for the integral of a nonnegative Borel function $f$ on $\RR_+$ with respect to this $\sigma$-finite measure. Furthermore, we can define
$$
\int_0^\infty f(s) \, da(s) = \lim_{T \to \infty} \int_0^T f(s) \, da(s) \in (-\infty, \infty)
$$
for any real Borel function $f$ on $\RR_+$ such that $\int_0^\infty |f(s)| \, |da(s)| < \infty$.


\subsection{Finite Variation Processes}
We now consider random variables and processes defined on a filtered probability space $(\Omega, \F, (\F_t), P)$.

\defnr{Finite Variation process}{finit_var_proc}{An adapted process $A = (A_t)_{t \geq 0}$ is called a \textit{finite variation process} if all its sample paths are finite variation functions on $\RR_+$. If in addition the sample paths are nondecreasing functions, the process $A$ is called an \textit{increasing process}.}
\textbf{Remark} In particular, $A_0 = 0$ and the sample paths of $A$ are continuous -- one can define finite variation processes with càdlàg sample paths, but in this book we consider \textbf{only the case of continuous sample paths}. Our special convention that the initial value of a finite variation process is 0 will be convenient for certain uniqueness statements.
If $A$ is a finite variation process, the process
\[ V_t = \int_0^t |dA_s| \]
is an increasing process. Indeed, it is clear that the sample paths of $V$ are nondecreasing functions (as well as continuous functions that vanish at $t = 0$). The fact that $V_t$ is an $\F_t$-measurable random variable can be deduced from the second part of Proposition \ref{prop:characterization_variation_func}. Writing $A_t = \frac{1}{2}(V_t + A_t) - \frac{1}{2}(V_t - A_t)$ shows that any finite variation process can be written as the difference of two increasing processes (the converse is obvious).

\proppr{finite_variation_proc_integral}{Let $A$ be a finite variation process, and let $H$ be a progressive process such that
$$ \forall t \geq 0, \, \forall \omega \in \Omega, \int_0^t |H_s(\omega)| \, |dA_s(\omega)| < \infty. $$
Then the process $H \cdot A = ((H \cdot A)_t)_{t \geq 0}$ defined by
$$ (H \cdot A)_t = \int_0^t H_s \, dA_s $$
is also a finite variation process.
}{
By the observations preceding the statement of Proposition \ref{prop:characterization_variation_func}, we know that the sample paths of $H \cdot A$ are finite variation functions. It remains to verify that the process $H \cdot A$ is adapted. To this end, it is enough to check that, if $t > 0$ is fixed, if $h : \Omega \times [0, t] \longrightarrow \RR$ is measurable for the $\sigma$-field $\F_t \otimes \B([0, t])$, and if $\int_0^t |h(\omega, s)| \, |dA_s(\omega)| < \infty$ for every $\omega$, then the variable $\int_0^t h(\omega, s) \, dA_s(\omega)$ is $\F_t$-measurable.

If $h(\omega, s) = \1_{(u, v]}(s) \1_\Gamma(\omega)$ with $(u, v] \subset [0, t]$ and $\Gamma \in \F_t$, the result is immediate since $\int_0^t h(\omega, s) \, dA_s(\omega) = \1_\Gamma(\omega) (A_v(\omega) - A_u(\omega))$ in that case. A monotone class argument \ref{thm:monotone_class_theorem} then gives the case $h = \1_G$, $G \in \F_t \otimes \B([0, t])$. Finally, in the general case, we observe that we can write $h$ as a pointwise limit of a sequence of simple functions (i.e., finite linear combinations of indicator functions of measurable sets) $h_n$ such that $|h_n| \leq |h|$ for every $n$, and that we then have $\int_0^t h_n(\omega, s) \, dA_s(\omega) \longrightarrow \int_0^t h(\omega, s) \, dA_s(\omega)$ by dominated convergence, for every $\omega \in \Omega$. \hfill $\square$}

\textbf{Remarks}
\begin{enumerate}
    \item[(i)] It happens frequently that instead of the assumption of the proposition we have the weaker assumption
    \[ \text{a.s. } \forall t \geq 0, \int_0^t |H_s(\omega)| \, |dA_s(\omega)| < \infty. \]
    If the filtration is complete, we can still define $H \cdot A$ as a finite variation process under this weaker assumption. We replace $H$ by the process $H'$ defined by
    \[
    H'_t(\omega) = \begin{cases}
    H_t(\omega) & \text{if } \int_0^n |H_s(\omega)| \, |dA_s(\omega)| < \infty, \, \forall n, \\
    0 & \text{otherwise.}
    \end{cases}
    \]
    Thanks to the fact that the filtration is complete, the process $H'$ is still progressive, which allows us to define $H \cdot A = H' \cdot A$. We will use this extension of Proposition \ref{prop:finite_variation_proc_integral} implicitly in what follows.
    \item[(ii)] Under appropriate assumptions (if $H$ and $K$ are progressive and $\int_0^t |H_s| \, |dA_s| < \infty$, $\int_0^t |H_s K_s| \, |dA_s| < \infty$ for every $t \geq 0$), we have the \textit{associativity property}
    \begin{equation*}
    K \cdot (H \cdot A) = (KH) \cdot A.
    \end{equation*}
    This indeed follows from the analogous deterministic result saying informally that $k(s) (h(s) \, \mu(ds)) = (k(s)h(s)) \, \mu(ds)$ if $h(s)$ and $k(s)h(s)$ are integrable with respect to the signed measure $\mu$ on $[0, t]$.
\end{enumerate}

An important special case of the proposition is the case where $A_t = t$. If $H$ is a progressive process such that
$$
\forall t \geq 0, \, \forall \omega \in \Omega, \int_0^t |H_s(\omega)| \, ds < \infty,
$$
the process $\int_0^t H_s \, ds$ is a finite variation process.

\section{Continuous Local Martingales}
We consider again a filtered probability space $(\Omega, \mathcal{F}, (\mathcal{F}_t), P)$. If $T$ is a stopping time, and if $X = (X_t)_{t \ge 0}$ is an adapted process with continuous sample paths, we will write $X^T$ for process $X$ stopped at $T$, defined by $X_t^T = X_{t \wedge T}$ for every $t \ge 0$. It is useful to observe that, if $S$ is another stopping time,
\[
(X^T)^S = (X^S)^T = X^{S \wedge T}.
\]
\subsection{Definition and Basic Properties of Continuous Local Martingales}
\defnr{continuous local martingale}{local_martingale}{
    An adapted process $M = (M_t)_{t \ge 0}$ with continuous sample paths and such that $M_0 = 0$ a.s. is called a \textit{continuous local martingale} if there exists a nondecreasing sequence $(T_n)_{n \ge 0}$ of stopping times such that $T_n \uparrow \infty$ (i.e. $T_n(\omega) \uparrow \infty$ for every $\omega$) and, for every $n$, the stopped process $M^{T_n}$ is a uniformly integrable martingale. 
    }
More generally, when we do not assume that $M_0 = 0$ a.s., we say that $M$ is a continuous local martingale if the process $N_t = M_t - M_0$ is a continuous local martingale. In all cases, we say that the sequence of stopping times $(T_n)$ \textit{reduces} $M$ if $T_n \uparrow \infty$ and, for every $n$, the stopped process $M^{T_n}$ is a uniformly integrable martingale.
\textbf{Remarks}
\begin{enumerate}
    \item We do not require in the definition of a continuous local martingale that the variables $M_t$ are in $L^1$ (compare with the definition of martingales). In particular, the variable $M_0$ may be any $\mathcal{F}_0$-measurable random variable.
    \item Any martingale with continuous sample paths is a continuous local martingale (see property (a) below) but the converse is false, and for this reason we will sometimes speak of \textit{true martingales} to emphasize the difference with local martingales.
    \item One can define a notion of local martingale with càdlàg sample paths. In this course, however, we consider only continuous local martingales.
\end{enumerate}
The following properties are easily established.
\propr{properties_local_martingale}{
\begin{enumerate}
\item A martingale with continuous sample paths is a continuous local martingale, and the sequence $T_n = n$ reduces $M$.

\item In the definition of a continuous local martingale starting from $0$, one can replace \text{uniformly integrable martingale} by \textit{martingale} (indeed, one can then observe that $M^{T_n \wedge n}$ is uniformly integrable, and we still have $T_n \wedge n \uparrow \infty$).

\item If $M$ is a continuous local martingale, then, for every stopping time $T$, $M^T$ is a continuous local martingale (this follows from Corollary of Optional Stopping Theorem \ref{thm:optional_stopping_martingale}).

\item If $(T_n)$ reduces $M$ and if $(S_n)$ is a sequence of stopping times such that $S_n \uparrow \infty$, then the sequence $(T_n \wedge S_n)$ also reduces $M$ (use same Corollary again).

\item The space of all continuous local martingales is a vector space (to check stability under addition, note that if $M$ and $M'$ are two continuous local martingales such that $M_0 = 0$ and $M'_0 = 0$, if the sequence $(T_n)$ reduces $M$ and if the sequence $(T'_n)$ reduces $M'$, property (d) shows that the sequence $T_n \wedge T'_n$ reduces $M + M'$).
\end{enumerate}}
The next proposition gives three other useful properties of local martingales.
\propr{other_properties_local_martingale}{
\begin{enumerate}
    \item A \textit{nonnegative continuous local martingale} $M$ such that $M_0 \in L^1$ is a \textit{supermartingale}.
    \item A \textit{continuous local martingale} $M$ such that there exists a random variable $Z \in L^1$ with $|M_t| \le Z$ for every $t \ge 0$ (in particular a \textit{bounded continuous local martingale}) is a \textit{uniformly integrable martingale}.
    \item If $M$ is a \textit{continuous local martingale} and $M_0 = 0$ (or more generally $M_0 \in L^1$), the sequence of stopping times
    \[T_n = \inf\{t \ge 0 : |M_t| \ge n\}\]
    \textit{reduces} $M$.
\end{enumerate}
}

\thmpr{Total variation of Local Martingales}{tot_var_local_martingal}{
Let $M$ be a continuous local martingale. Assume that $M$ is also a finite variation process (in particular $M_0 = 0$). Then $M_t = 0$ for every $t \ge 0$, a.s.
}{
Set
\[\tau_n = \inf\left\{t \ge 0 : \int_0^t |dM_s| \ge n\right\}\] 
for every integer $n \ge 0$. By Definition \ref{defn:hitting_time}, $\tau_n$ is a stopping time (recall that $\int_0^t |dM_s|$ is an increasing process if $M$ is a finite variation process).

Fix $n \ge 0$ and set $N = M^{\tau_n}$. Note that, for every $t \ge 0$,
\[ |N_t| = |M_{t \wedge \tau_n}| \le \int_0^t \wedge \tau_n |dM_s| \le n. \]
By Proposition \ref{prop:other_properties_local_martingale}, $N$ is a (bounded) martingale. Let $t > 0$ and let $0 = t_0 < t_1 < \dots < t_p = t$ be any subdivision of $[0, t]$. Then, from Proposition \ref{prop:finite_quad_var}, we have
\[
\begin{aligned}
E[N_t^2] &= \sum_{i=1}^p E[(N_{t_i} - N_{t_{i-1}})^2] \\
&\le E\left[\left(\sup_{1 \le i \le p} |N_{t_i} - N_{t_{i-1}}|\right) \sum_{i=1}^p |N_{t_i} - N_{t_{i-1}}|\right] \\
&\le n \, E\left[\sup_{1 \le i \le p} |N_{t_i} - N_{t_{i-1}}|\right]
\end{aligned}
\]
noting that $\int_0^t |dN_s| \le n$ by the definition of $\tau_n$, and using Proposition \ref{prop:characterization_variation_func}.

We now apply the preceding bound to a sequence $0 = t_0^k < t_1^k < \dots < t_{p_k}^k = t$ of subdivisions of $[0, t]$ whose mesh tends to 0. Using the continuity of sample paths, and the fact that $N$ is bounded (to justify dominated convergence), we get
\[ \lim_{k \to \infty} E\left[\sup_{1 \le i \le p_k} |N_{t_i^k} - N_{t_{i-1}^k}|\right] = 0. \]
We then conclude that $E[N_t^2] = 0$, hence $M_{t \wedge \tau_n} = 0$ a.s. Letting $n$ tend to $\infty$, we get that $M_t = 0$ a.s. 
}
\subsection{The Quadratic Variation of a Continuous Local Martingale} \label{sec:quadratic-variation}

\thmpr{Quadratic variation existence}{quad_var_existence}{
Let $M = (M_t)_{t \geq 0}$ be a continuous local martingale. There exists an increasing process denoted by $(\langle M, M \rangle_t)_{t \geq 0}$, which is unique up to indistinguishability, such that $M_t^2 - \langle M, M \rangle_t$ is a continuous local martingale. Furthermore, for every fixed $t > 0$, if $0 = t_0^n < t_1^n < \dots < t_{p_n}^n = t$ is an increasing sequence of subdivisions of $[0, t]$ with mesh tending to $0$, we have :
\begin{align*}
\langle M, M \rangle_t = \lim_{n \to \infty} \sum_{i=1}^{p_n} (M_{t_i^n} - M_{t_{i-1}^n})^2 
\end{align*}
In probability. The process $\langle M, M \rangle$ is called the quadratic variation of $M$.
}{
    The proof can be found in \parencite[thm.~4.9]{Stochastic}
}
\textbf{Remarks:}
\begin{enumerate}
    \item We observe that the process $\langle M, M \rangle$ does not depend on the initial value $M_0$, but only on the increments of $M$: if $M_t = M_0 + N_t$, we have $\langle M, M \rangle = \langle N, N \rangle$. This is obvious from the second assertion of the theorem, and this will also be clear from the proof that follows.

    \item In the second assertion of the theorem, it is in fact not necessary to assume that the sequence of subdivisions is increasing.
\end{enumerate}
Now we present some properties of the quadratic variation : 
\proppr{quad_var_propt}{
Let $M$ be a continuous local martingale and let $T$ be a stopping time. 

\begin{enumerate}
    \item We have a.s. for every $t \geq 0$,
    \begin{align*}
    \langle M^T, M^T \rangle_t = \langle M, M \rangle_{t \wedge T}.
    \end{align*}
    \item If $M_0 = 0$. Then we have $\langle M, M \rangle = 0$ if and only if $M = 0$.
\end{enumerate}
}{
\begin{enumerate}
    \item This follows from the fact that $M_{t \wedge T}^2 - \langle M, M \rangle_{t \wedge T}$ is a continuous local martingale (cf. property (3) of continuous local martingales in Proposition \ref{prop:properties_local_martingale}).
    \item Suppose that $\langle M, M \rangle = 0$. Then $M_t^2$ is a nonnegative continuous local martingale and, by Proposition \ref{prop:other_properties_local_martingale} (1), $M_t^2$ is a supermartingale, hence $E[M_t^2] \leq E[M_0^2] = 0$, so that $M_t = 0$ for every $t$. The converse is obvious.
\end{enumerate}
}
The next theorem shows that properties of a continuous local martingale are closely related to those of its quadratic variation. If $A$ is an increasing process, $A_\infty$ denotes the increasing limit of $A_t$ as $t \to \infty$ (this limit always exists in $[0, \infty]$).

\thmpr{}{quad_var_equivalences}{
Let $M$ be a continuous local martingale with $M_0 \in L^2$.
\begin{enumerate}
\item \textit{The following are equivalent:}
\begin{itemize}
    \item[(a)] \textit{$M$ is a (true) martingale bounded in $L^2$.}
    \item[(b)] \textit{$E[\langle M, M \rangle_\infty] < \infty$.}
\end{itemize}
\textit{Furthermore, if these properties hold, the process $M_t^2 - \langle M, M \rangle_t$ is a uniformly integrable martingale, and in particular $E[M_\infty^2] = E[M_0^2] + E[\langle M, M \rangle_\infty]$.}

\item \textit{The following are equivalent:}
\begin{itemize}
    \item[(a)] \textit{$M$ is a (true) martingale and $M_t \in L^2$ for every $t \geq 0$.}
    \item[(b)] \textit{$E[\langle M, M \rangle_t] < \infty$ for every $t \geq 0$.}
\end{itemize}
\end{enumerate}
Furthermore, if these properties hold, the process $M_t^2 - \langle M, M \rangle_t$ is a martingale.
}{
Let's prove (1), Replacing $M$ by $M - M_0$, we may assume that $M_0 = 0$ in the proof. Let us first assume that $M$ is a martingale bounded in $L^2$. Doob's inequality in $L^2$ (Proposition \ref{thm:doob_supermartingale_inequalities} (4)) shows that, for every $T > 0$,
\begin{align*}
E\left[ \sup_{0 \leq t \leq T} M_t^2 \right] \leq 4 E[M_T^2].
\end{align*}
By letting $T$ go to $\infty$, we have
\begin{align*}
E\left[ \sup_{t \geq 0} M_t^2 \right] \leq 4 \sup_{t \geq 0} E[M_t^2] =: C < \infty.
\end{align*}

Set $S_n = \inf\{t \geq 0 : \langle M, M \rangle_t \geq n\}$. Then the continuous local martingale $M_{t \wedge S_n}^2 - \langle M, M \rangle_{t \wedge S_n}$ is dominated by the variable
\begin{align*}
\sup_{s \geq 0} M_s^2 + n,
\end{align*}
which is integrable. From Proposition \ref{prop:other_properties_local_martingale} (2), we get that this continuous local martingale is a uniformly integrable martingale, hence
\begin{align*}
E[\langle M, M \rangle_{t \wedge S_n}] = E[M_{t \wedge S_n}^2] \leq E\left[ \sup_{s \geq 0} M_s^2 \right] \leq C.
\end{align*}
By letting $n$, and then $t$ tend to infinity, and using monotone convergence, we get $E[\langle M, M \rangle_\infty] \leq C < \infty$.

Conversely, assume that $E[\langle M, M \rangle_\infty] < \infty$. Set $T_n = \inf\{t \geq 0 : |M_t| \geq n\}$. Then the continuous local martingale $M_{t \wedge T_n}^2 - \langle M, M \rangle_{t \wedge T_n}$ is dominated by the variable
\begin{align*}
n^2 + \langle M, M \rangle_\infty,
\end{align*}
which is integrable. From Proposition \ref{prop:other_properties_local_martingale} (2) again, this continuous local martingale is a uniformly integrable martingale, hence, for every $t \geq 0$,
\begin{align*}
E[M_{t \wedge T_n}^2] = E[\langle M, M \rangle_{t \wedge T_n}] \leq E[\langle M, M \rangle_\infty] =: C' < \infty.
\end{align*}

By letting $n \to \infty$ and using Fatou's lemma, we get $E[M_t^2] \leq C'$, so that the collection $(M_t)_{t \geq 0}$ is bounded in $L^2$. We have not yet verified that $(M_t)_{t \geq 0}$ is a martingale. However, the previous bound on $E[M_{t \wedge T_n}^2]$ shows that the sequence $(M_{t \wedge T_n})_{n \geq 1}$ is uniformly integrable, and therefore converges both a.s. and in $L^1$ to $M_t$, for every $t \geq 0$. Recalling that $M^{T_n}$ is a martingale (Proposition \ref{prop:other_properties_local_martingale} (3)), the $L^1$-convergence allows us to pass to the limit $n \to \infty$ in the martingale property $E[M_{t \wedge T_n} \mid \mathcal{F}_s] = M_{s \wedge T_n}$, for $0 \leq s < t$, and to get that $M$ is a martingale.

Finally, if properties (a) and (b) hold, the continuous local martingale $M^2 - \langle M, M \rangle$ is dominated by the integrable variable
\begin{align*}
\sup_{t \geq 0} M_t^2 + \langle M, M \rangle_\infty
\end{align*}
and is therefore (by Proposition \ref{prop:other_properties_local_martingale} (2)) a uniformly integrable martingale.
}
\textbf{Remark:}
\begin{enumerate}
    \item In property (a) of (1) (or of (2)), it is essential to suppose that $M$ is a martingale, and not only a continuous local martingale. Doob's inequality used in the previous proof is not valid in general for a continuous local martingale!
    \item It suffices to apply (1) to $(M_{t \wedge a})_{t \geq 0}$ for every choice of $a \geq 0$, to prove (2). 
\end{enumerate}
    

\lempr{Quadratic-variation integral formula for continuous integrands}{qv_continuous_integrand}{
Let $X = (X_s)_{s \ge 0}$ be a continuous process for which Theorem~\ref{thm:quad_var_existence} holds, so that the quadratic variation $\langle X, X \rangle$ exists as a continuous, nondecreasing, adapted process. Fix $t > 0$ and let $(\pi_n)_{n \ge 1}$,
$$
\pi_n : \quad 0 = t_0^n < t_1^n < \cdots < t_{p_n}^n = t,
$$
be a sequence of subdivisions of $[0,t]$ whose mesh $|\pi_n| := \max_i (t_{i+1}^n - t_i^n)$
tends to $0$ as $n \to \infty$.
Let $H = (H_s)_{0 \le s \le t}$ be a process, jointly measurable in $(s,\omega)$, whose
sample paths $s \mapsto H_s(\omega)$ are almost surely continuous. Then
$$
\sum_{i=0}^{p_n - 1} H_{t_i^n} \left( X_{t_{i+1}^n} - X_{t_i^n} \right)^2
\;\xrightarrow[n \to \infty]{\mathbb{P}}\;
\int_0^t H_s \, d\langle X, X \rangle_s .
$$
}{
Since convergence in probability is metrizable, it suffices to show that every
subsequence of $(n)$ has a further subsequence along which the stated convergence holds almost surely.
\begin{enumerate}
\item \textbf{Step 1 (Random measures).}
For each $n$, define the (a.s. finite, purely atomic, positive) random measure on $[0,t]$
$$
\mu_n(dr) \;:=\; \sum_{i=0}^{p_n - 1} \left( X_{t_{i+1}^n} - X_{t_i^n} \right)^2
\delta_{t_i^n}(dr).
$$
By construction,
$$
\int_{[0,t]} H_s \, \mu_n(ds) \;=\; \sum_{i=0}^{p_n - 1} H_{t_i^n}
\left( X_{t_{i+1}^n} - X_{t_i^n} \right)^2. \quad [1]
$$

\item \textbf{Step 2 (Convergence of cumulative masses via Proposition~\ref{thm:quad_var_existence}).}
Fix $r \in [0,t]$ and let $k = k_n(r) := \max\{\, i : t_i^n \le r \,\}$. Applying
Proposition~\ref{thm:quad_var_existence} (with $Y = X$) to the subdivision
$t_0^n < t_1^n < \cdots < t_k^n < r$ of $[0,r]$ — whose mesh also tends to $0$ — gives
$$
\sum_{i=0}^{k-1} \left( X_{t_{i+1}^n} - X_{t_i^n} \right)^2
+ \left( X_r - X_{t_k^n} \right)^2
\;\xrightarrow[n \to \infty]{\mathbb{P}}\; \langle X, X \rangle_r .
$$
On the other hand, directly from the definition of $\mu_n$,
$$
\mu_n([0,r]) \;=\; \sum_{i=0}^{k-1} \left( X_{t_{i+1}^n} - X_{t_i^n} \right)^2
+ \left( X_{t_{k+1}^n} - X_{t_k^n} \right)^2 .
$$
The two displays differ only in their last term, and
$$
\left( X_{t_{k+1}^n} - X_{t_k^n} \right)^2 - \left( X_r - X_{t_k^n} \right)^2
\;\xrightarrow[n \to \infty]{\text{a.s.}}\; 0,
$$
since $t_{k+1}^n - t_k^n \to 0$ and $r - t_k^n \to 0$ as $|\pi_n| \to 0$, and $X$ has
(a.s. uniformly) continuous paths on $[0,t]$. Hence
$$
\mu_n([0,r]) \;\xrightarrow[n \to \infty]{\mathbb{P}}\; \langle X, X \rangle_r
\qquad \text{for every } r \in [0,t]. \quad [1]
$$

\item \textbf{Step 3 (Diagonal extraction).}
Let $D := \bigcup_{n \ge 1} \{ t_0^n, \dots, t_{p_n}^n \}$, a countable dense subset of
$[0,t]$ containing $t$. Enumerate $D = \{r_1, r_2, \dots\}$. By $(**)$ and a diagonal
argument, there is a subsequence $(n_k)_{k \ge 1}$ along which
$$
\mu_{n_k}([0, r_j]) \;\xrightarrow[k \to \infty]{\text{a.s.}}\; \langle X, X \rangle_{r_j}
\qquad \text{simultaneously for every } j \ge 1.
$$

\item \textbf{Step 4 (Weak convergence of $\mu_{n_k}$).}
Let $\Omega_0$ be the a.s. event on which: (i) the convergence of Step~3 holds at every
$r \in D$; (ii) $s \mapsto \langle X,X\rangle_s(\omega)$ is continuous and nondecreasing;
(iii) $s \mapsto H_s(\omega)$ is continuous. Fix $\omega \in \Omega_0$. The nondecreasing
functions $r \mapsto \mu_{n_k}([0,r])(\omega)$ converge on the dense set $D$ to the
continuous nondecreasing function $r \mapsto \langle X,X \rangle_r(\omega)$; a standard
squeeze argument using monotonicity and density then extends this to convergence at
\emph{every} $r \in [0,t]$, including $r = t$ (total mass). Convergence of these
"distribution functions" at every point, together with convergence of total mass, implies
$$
\mu_{n_k}(\omega) \;\xrightarrow[k \to \infty]{w}\;
\mathbf{1}_{[0,t]}(r) \, d\langle X, X \rangle_r(\omega)
$$
weakly as finite positive measures on $[0,t]$.

\item \textbf{Step 5 (Passing to the limit against $H$).}
Since $H_\cdot(\omega)$ is continuous, hence bounded, on the compact interval $[0,t]$,
weak convergence yields
$$
\int_{[0,t]} H_s(\omega) \, \mu_{n_k}(ds)(\omega)
\;\xrightarrow[k \to \infty]{}\;
\int_0^t H_s(\omega) \, d\langle X, X \rangle_s(\omega),
\qquad \text{for every } \omega \in \Omega_0.
$$
By $(*)$, this is exactly
$$
\sum_{i=0}^{p_{n_k}-1} H_{t_i^{n_k}}
\left( X_{t_{i+1}^{n_k}} - X_{t_i^{n_k}} \right)^2
\;\xrightarrow[k \to \infty]{\text{a.s.}}\;
\int_0^t H_s \, d\langle X, X \rangle_s .
$$

\item \textbf{Step 6 (Conclusion).}
Every subsequence of $(n)$ admits a further subsequence along which the convergence
holds almost surely; therefore the original sequence converges in probability, which
proves the lemma.
\end{enumerate}
}

\subsection{The Bracket of Two Continuous Local Martingales}
\defnr{Bracket of two continuous local martingales}{martingale_bracket}{
If $M$ and $N$ are two continuous local martingales, the \textit{bracket} $\langle M, N \rangle$ is the finite variation process defined by setting, for every $t \geq 0$,
\begin{align*}
\langle M, N \rangle_t = \frac{1}{2} (\langle M + N, M + N \rangle_t - \langle M, M \rangle_t - \langle N, N \rangle_t).
\end{align*}
}
Let us state a few easy properties of the bracket.
\propr{Bracket_properties}{
\begin{enumerate}
    \item \textit{$\langle M, N \rangle$ is the unique (up to indistinguishability) finite variation process such that $M_t N_t - \langle M, N \rangle_t$ is a continuous local martingale.}
    \item  \textit{The mapping $(M, N) \mapsto \langle M, N \rangle$ is bilinear and symmetric.}
    \item  \textit{If $0 = t_0^n < t_1^n < \dots < t_{p_n}^n = t$ is an increasing sequence of subdivisions of $[0, t]$ with mesh tending to $0$, we have}
    \begin{align*}
    \lim_{n \to \infty} \sum_{i=1}^{p_n} (M_{t_i^n} - M_{t_{i-1}^n})(N_{t_i^n} - N_{t_{i-1}^n}) = \langle M, N \rangle_t
    \end{align*}
    \textit{in probability.}
    \item  \textit{For every stopping time $T$, $\langle M^T, N^T \rangle_t = \langle M^T, N \rangle_t = \langle M, N \rangle_{t \wedge T}$.}
    \item \textit{If $M$ and $N$ are two martingales (with continuous sample paths) bounded in $L^2$, $M_t N_t - \langle M, N \rangle_t$ is a uniformly integrable martingale. Consequently, $\langle M, N \rangle_\infty$ is well defined as the almost sure limit of $\langle M, N \rangle_t$ as $t \to \infty$, is integrable, and satisfies}
    \begin{align*}
    \EE[M_\infty N_\infty] = \EE[M_0 N_0] + \EE[\langle M, N \rangle_\infty].
    \end{align*}
\end{enumerate}
}
\textbf{Remark:} A consequence of (iv) is the fact that $M^T(N-N^T)$ is a continuous local martingale, which is not so easy to prove directly.

\defnr{Orthogonal continuous local martingales}{orthogonal_local_martingales}{
Two continuous local martingales $M$ and $N$ are said to be orthogonal if $\langle M, N \rangle = 0$, which holds if and only if $MN$ is a continuous local martingale.}
If $M$ and $N$ are two orthogonal martingales bounded in $L^2$, we have $E[M_t N_t] = E[M_0 N_0]$, and even $E[M_S N_S] = E[M_0 N_0]$ for any stopping time $S$. This follows from Theorem \ref{thm:optional_stopping_martingale}, using property (5) of Proposition \ref{prop:Bracket_properties}.

\thmpr{Kunita-Watanabe inequality}{kunita_watanabi_inequality}{
    Let $M$ and $N$ be two continuous local martingales and let $H$ and $K$ be two measurable processes. Then, a.s.,
    \begin{align*}
    \int_0^\infty |H_s| |K_s| |{\rm d}\langle M, N \rangle_s| \leq \left( \int_0^\infty H_s^2 {\rm d}\langle M, M \rangle_s \right)^{1/2} \left( \int_0^\infty K_s^2 {\rm d}\langle N, N \rangle_s \right)^{1/2}.
    \end{align*}
}{
Only in this proof, we use the special notation $\langle M, N \rangle_s^t = \langle M, N \rangle_t - \langle M, N \rangle_s$ for $0 \leq s \leq t$. The first step of the proof is to observe that we have a.s. for every choice of the rationals $s < t$ (and also by continuity for every reals $s < t$),
\begin{align*}
|\langle M, N \rangle_s^t| \leq \sqrt{\langle M, M \rangle_s^t} \sqrt{\langle N, N \rangle_s^t}.
\end{align*}

Indeed, this follows from the approximations of $\langle M, M \rangle$ and $\langle M, N \rangle$ given in Theorem \ref{thm:quad_var_existence} and in Proposition \ref{prop:Bracket_properties} respectively (note that these approximations are easily extended to the increments of $\langle M, M \rangle$ and $\langle M, N \rangle$), together with the Cauchy-Schwarz inequality. From now on, we fix $\omega$ such that the inequality of the last display holds for every $s < t$, and we argue with this value of $\omega$ (the remaining part of the argument is \textit{deterministic}).

We then observe that we also have, for every $0 \leq s \leq t$,
\begin{align*}
\int_s^t |{\rm d}\langle M, N \rangle_u| \leq \sqrt{\langle M, M \rangle_s^t} \sqrt{\langle N, N \rangle_s^t}. \tag{*}
\end{align*}

Indeed, we use Proposition \ref{prop:characterization_variation_func}, noting that, for any subdivision $s = t_0 < t_1 < \dots < t_p = t$, we can bound
\begin{align*}
\sum_{i=1}^p |\langle M, N \rangle_{t_{i-1}}^{t_i}| &\leq \sum_{i=1}^p \sqrt{\langle M, M \rangle_{t_{i-1}}^{t_i}} \sqrt{\langle N, N \rangle_{t_{i-1}}^{t_i}} \\
&\leq \left( \sum_{i=1}^p \langle M, M \rangle_{t_{i-1}}^{t_i} \right)^{1/2} \left( \sum_{i=1}^p \langle N, N \rangle_{t_{i-1}}^{t_i} \right)^{1/2} \\
&= \sqrt{\langle M, M \rangle_s^t} \sqrt{\langle N, N \rangle_s^t}.
\end{align*}

We then get that, for every bounded Borel subset $A$ of $\mathbb{R}_+$,
\begin{align*}
\int_A |{\rm d}\langle M, N \rangle_u| \leq \sqrt{\int_A {\rm d}\langle M, M \rangle_u} \sqrt{\int_A {\rm d}\langle N, N \rangle_u}.
\end{align*}

When $A = [s, t]$, this is the bound (*). If $A$ is a finite union of intervals, this follows from (*) and another application of the Cauchy-Schwarz inequality. A monotone class argument shows that the inequality of the last display remains valid for any bounded Borel set $A$ (Theorem \ref{thm:monotone_class_theorem}).

Next let $h = \sum_{i=1}^p \lambda_i \mathbf{1}_{A_i}$ and $k = \sum_{i=1}^p \mu_i \mathbf{1}_{A_i}$ be two nonnegative simple functions on $\mathbb{R}_+$ with bounded support contained in $[0, K]$, for some $K > 0$. Here $A_1, \dots, A_p$ is a measurable partition of $[0, K]$, and $\lambda_1, \dots, \lambda_p, \mu_1, \dots, \mu_p$ are reals (we can always assume that $h$ and $k$ are expressed in terms of the same partition). Then,
\begin{align*}
\int h(s)k(s)|{\rm d}\langle M, N \rangle_s| &= \sum_{i=1}^p \lambda_i \mu_i \int_{A_i} |{\rm d}\langle M, N \rangle_s| \\
&\leq \sum_{i=1}^p \lambda_i \mu_i \sqrt{\int_{A_i} {\rm d}\langle M, M \rangle_s} \sqrt{\int_{A_i} {\rm d}\langle N, N \rangle_s} \\
&\leq \left( \sum_{i=1}^p \lambda_i^2 \int_{A_i} {\rm d}\langle M, M \rangle_s \right)^{1/2} \left( \sum_{i=1}^p \mu_i^2 \int_{A_i} {\rm d}\langle N, N \rangle_s \right)^{1/2} \\
&= \left( \int h(s)^2 {\rm d}\langle M, M \rangle_s \right)^{1/2} \left( \int k(s)^2 {\rm d}\langle N, N \rangle_s \right)^{1/2},
\end{align*}
which gives the desired inequality for simple functions. Since every nonnegative Borel function is a monotone increasing limit of simple functions with bounded support, an application of the monotone convergence theorem completes the proof. 
}


\subsection{The Hilbert Space \texorpdfstring{$\mathbb{H}^2$}{H2}}

We write $\mathbb{H}^2$ for the space of all continuous martingales $M$ which
are bounded in $L^2$ and such that $M_0=0$, with the usual convention that two
indistinguishable processes are identified. Equivalently, $M\in\mathbb{H}^2$
if and only if $M$ is a continuous local martingale such that $M_0=0$ and
\[
\EE[\langle M,M\rangle_\infty]<\infty;
\]
see Theorem \ref{thm:quad_var_equivalences}. By Theorem
\ref{thm:Doob_UI_convergence}, every $M\in\mathbb{H}^2$ admits an almost sure
and $L^2$ limit $M_\infty\in L^2$, and
\[
M_t=\EE[M_\infty\mid\F_t],\qquad t\geq 0.
\]

If $M,N\in\mathbb{H}^2$, Proposition \ref{prop:Bracket_properties} (v) shows
that $\langle M,N\rangle_\infty$ is well defined and integrable. We may
therefore define a symmetric bilinear form on $\mathbb{H}^2$ by
\[
(M,N)_{\mathbb{H}^2}
=\EE[\langle M,N\rangle_\infty]
=\EE[M_\infty N_\infty].
\]
The second equality follows again from Proposition
\ref{prop:Bracket_properties} (v), since $M_0=N_0=0$. In particular,
$(M,M)_{\mathbb{H}^2}=0$ if and only if $M=0$. The associated norm is
\[
\|M\|_{\mathbb{H}^2}
=\bigl((M,M)_{\mathbb{H}^2}\bigr)^{1/2}
=\left(\EE[\langle M,M\rangle_\infty]\right)^{1/2}
=\left(\EE[M_\infty^2]\right)^{1/2}.
\]

\proppr{hilbert_structure_H2}{
The space $\mathbb{H}^2$, equipped with the scalar product
$(\cdot,\cdot)_{\mathbb{H}^2}$, is a Hilbert space.
}{
Let $(M^n)_{n\geq 1}$ be a Cauchy sequence in $\mathbb{H}^2$. Then
$(M^n_\infty)_{n\geq 1}$ is Cauchy in $L^2$, hence converges in $L^2$ to some
$X\in L^2$. For each $t\geq 0$, the sequence $(M^n_t)_{n\geq 1}$ converges in
$L^2$ to $M_t:=\EE[X\mid\F_t]$. Since
$\EE[M^n_\infty\mid\F_0]=M^n_0=0$ for every $n$, we also have
$M_0=\EE[X\mid\F_0]=0$. Thus $(M_t)_{t\geq 0}$ is a martingale bounded in
$L^2$.

Moreover, Doob's $L^2$ inequality gives, for every $n,m\geq 1$,
\[
\EE\left[\sup_{t\geq 0}|M^n_t-M^m_t|^2\right]
\leq 4\,\EE\left[|M^n_\infty-M^m_\infty|^2\right].
\]
By extracting a subsequence, we obtain almost sure uniform convergence of
$M^n$ to $M$ on $\RR_+$. Since each $M^n$ has continuous sample paths, $M$
has a continuous modification. Hence $M\in\mathbb{H}^2$. Finally,
\[
\|M^n-M\|_{\mathbb{H}^2}^2
=\EE\left[|M^n_\infty-X|^2\right]\longrightarrow 0,
\]
which proves that $\mathbb{H}^2$ is complete.
}

\section{Continuous Semimartingales}
We now introduce the class of processes for which we will develop the theory of stochastic integrals.

\defnr{Continuous semimartingale}{continuous_semimartingale}{
A process $X = (X_t)_{t \ge 0}$ is a \textit{continuous semimartingale} if it can be written in the form
\[X_t = M_t + A_t,\]
where $M$ is a continuous local martingale and $A$ is a finite variation process.
}

The decomposition $X = M + A$ is then unique up to indistinguishability thanks to Theorem \ref{thm:tot_var_local_martingal}. We say that this is the \textit{canonical decomposition} of $X$.

By construction, continuous semimartingales have continuous sample paths. It is possible to define a notion of semimartingale with càdlàg sample paths, but in this book, we will only deal with \textit{continuous} semimartingales, and for this reason we sometimes omit the word continuous.

\defnr{Bracket of semimartingales}{bracket_semimartingales}{
Let $X = M + A$ and $Y = M' + A'$ be the canonical decompositions of two continuous semimartingales $X$ and $Y$. The \textit{bracket} $\langle X, Y \rangle$ is the finite variation process defined by
\[\langle X, Y \rangle_t = \langle M, M' \rangle_t.\]
In particular, we have $\langle X, X \rangle_t = \langle M, M \rangle_t$.
}

\proppr{prop_bracket_limit}{
Let $0 = t_0^n < t_1^n < \dots < t_{p_n}^n = t$ be an increasing sequence of subdivisions of $[0, t]$ whose mesh tends to 0. Then,
\[
\lim_{n \to \infty} \sum_{i=1}^{p_n} (X_{t_i^n} - X_{t_{i-1}^n})(Y_{t_i^n} - Y_{t_{i-1}^n}) = \langle X, Y \rangle_t
\]
in probability.
}{
We treat the case where $X = Y$ and leave the general case to the reader. We have
\[
\sum_{i=1}^{p_n} (X_{t_i^n} - X_{t_{i-1}^n})^2 = \sum_{i=1}^{p_n} (M_{t_i^n} - M_{t_{i-1}^n})^2 + \sum_{i=1}^{p_n} (A_{t_i^n} - A_{t_{i-1}^n})^2 + 2 \sum_{i=1}^{p_n} (M_{t_i^n} - M_{t_{i-1}^n})(A_{t_i^n} - A_{t_{i-1}^n}).
\]
By Theorem \ref{thm:quad_var_existence},
\[
\lim_{n \to \infty} \sum_{i=1}^{p_n} (M_{t_i^n} - M_{t_{i-1}^n})^2 = \langle M, M \rangle_t = \langle X, X \rangle_t,
\]
in probability. On the other hand,
\[
\begin{aligned}
\sum_{i=1}^{p_n} (A_{t_i^n} - A_{t_{i-1}^n})^2 &\le \left( \sup_{1 \le i \le p_n} |A_{t_i^n} - A_{t_{i-1}^n}| \right) \sum_{i=1}^{p_n} |A_{t_i^n} - A_{t_{i-1}^n}| \\
&\le \left( \int_0^t |dA_s| \right) \sup_{1 \le i \le p_n} |A_{t_i^n} - A_{t_{i-1}^n}|,
\end{aligned}
\]
which tends to 0 a.s. when $n \to \infty$ by the continuity of sample paths of $A$. The same argument shows that
\[
\left| \sum_{i=1}^{p_n} (A_{t_i^n} - A_{t_{i-1}^n})(M_{t_i^n} - M_{t_{i-1}^n}) \right| \le \left( \int_0^t |dA_s| \right) \sup_{1 \le i \le p_n} |M_{t_i^n} - M_{t_{i-1}^n}|
\]
tends to 0 a.s. 
}


\end{document}
